% ---------------------------------------------------------------------------
% Capítulo de desenvolvimento do protótipo.
% Fontes: repositórios BACKEND_TCC2_MATHEUS_PERILLO e FRONTEND_TCC2_MATHEUS_PERILLO
% (código, README, docker-compose.yml, testes, docs/desempenho-grafo.md),
% docs/mvp-escopo.md, docs/mvp-diario.md, docs/backlog/, docs/sistema/.
% Tabelas e figuras de dados geradas por docs/sistema/*.py; diagramas em
% editaveis/diagramas/ (fontes de cada elemento em docs/sistema/diagramas-fontes.md).
% ---------------------------------------------------------------------------
\begingroup
\setlength{\emergencystretch}{3em}

\chapter{Desenvolvimento do protótipo}
\label{cap:desenvolvimento}

Este capítulo descreve o protótipo do buscador: a arquitetura, os casos de
uso, o modelo de dados, os fluxos principais, a interface de programação, as
decisões de arquitetura, o modelo de ameaças, a engenharia de testes e o processo de
desenvolvimento com IA. O protótipo implementa a maior parte das histórias
dos três horizontes de contratações públicas do \textit{backlog}. Das 18
histórias do escopo de lavagem de dinheiro (H40 a H57), 7 estão feitas, 5
estão em parte, 5 aguardam decisão do orientador e 1 ficou fora desta
entrega. A situação de cada uma está na Tabela~\ref{tab:backlog}. O código
está em dois repositórios, sob a licença
MIT: um para o
\textit{back-end} e outro para o \textit{front-end}. Os testes automatizados
e a demonstração usam só dados sintéticos. Os conectores também foram
conferidos contra as fontes reais, num teste de fumaça
(Seção~\ref{sec:desenv-verificacao}).

% ===========================================================================
\section{Visão geral e arquitetura}
\label{sec:desenv-arquitetura}

O protótipo segue o fluxo de uma investigação sobre contratações públicas:
o investigador abre um caso com finalidade e prazo de retenção, coleta dados
abertos de várias fontes, o sistema une as entidades e monta o grafo de
vínculos, as tipologias geram alertas e o investigador confirma ou descarta
cada alerta. Toda captura entra na cadeia de custódia, e toda ação entra na
trilha de auditoria.

\textbf{Tecnologias.} O \textit{back-end} usa Python 3.12, FastAPI,
SQLAlchemy~2 e PostgreSQL (SQLite nos testes). O \textit{front-end} usa
React, TypeScript, Vite, Cytoscape.js para o grafo e o Padrão Digital de
Governo (\texttt{@govbr-ds/core}) sem a marca gov.br. A escolha buscou o
menor número de peças que entregasse a fatia toda, e é decisão do autor
(Tabela~\ref{tab:adr}).

Os diagramas deste capítulo usam a notação da UML~2, padronizada na ISO/IEC
19505-2 \cite{iso19505-2}. Cada elemento de cada diagrama foi conferido no
código ou na configuração (\texttt{docs/sistema/diagramas-fontes.md}). A Figura~\ref{fig:uml-componentes} mostra os módulos e as dependências entre
eles, e a Figura~\ref{fig:uml-implantacao}, a implantação em contêineres.

% Diagrama de componentes (UML 2) do buscador OSINT. Elementos e fontes no código:
% docs/sistema/diagramas-fontes.md, seção 2.
\begin{figure}[htbp]
\centering
\caption{Diagrama de componentes do buscador OSINT}
\label{fig:uml-componentes}
\resizebox{\textwidth}{!}{%
\begin{tikzpicture}[
  font=\small,
  comp/.style={draw, fill=white, minimum width=3.5cm, minimum height=1.15cm, align=center,
    inner sep=3pt},
  dep/.style={-{Straight Barb[length=6pt]}, dashed},
  sub/.style={draw, rounded corners=2pt, inner sep=0.35cm},
  titulo/.style={anchor=north west, font=\small\bfseries, inner sep=4pt},
  rot/.style={font=\scriptsize, fill=white, inner sep=1pt, align=center},
  icone/.pic={
    \draw[fill=white] (0,0) rectangle (0.26,0.32);
    \draw[fill=white] (-0.06,0.05) rectangle (0.07,0.12);
    \draw[fill=white] (-0.06,0.19) rectangle (0.07,0.26);}
]
\newcommand{\cmp}[2]{{\scriptsize«componente»}\\\textbf{#1}\\{\scriptsize\texttt{#2}}}
% ======================= front-end =======================
\node[comp] (telas) at (2.6,13.8)  {\cmp{telas}{src/telas/}};
\node[comp] (cmps)  at (7.8,13.8)  {\cmp{componentes}{src/componentes/}};
\node[comp] (sess)  at (13.6,13.8) {\cmp{sessão}{src/sessao/}};
\node[comp] (cli)   at (7.8,11.4)  {\cmp{cliente da API}{src/api/clienteHttp.ts}};
\node[comp] (sim)   at (15.2,11.4) {\cmp{cliente simulado}{src/api/clienteMock.ts}};
\draw[dep] (telas) -- (cmps);
\draw[dep] (cmps) -- (sess);
\draw[dep] (telas.north) -- ++(0,0.5) -| (sess.north);
\draw[dep] (telas.south east) -- (cli.north west);
\draw[dep] (cmps) -- (cli);
\draw[dep] (sess.south west) -- ([xshift=0.8cm]cli.north);
\draw[dep] (cli) -- node[rot, above=3pt]{import dinâmico,\\só com \texttt{VITE\_MOCK=1}} (sim);
\node[sub, fit={(telas)(sess)(sim) ($(telas.north)+(0,1.2)$)}] (FE) {};
\node[titulo] at (FE.north west) {Front-end (React, TypeScript), servido pelo nginx};
% ======================= interface =======================
\coordinate (bola) at (7.8,9.0);
\draw (bola) circle (0.17);
\draw ([yshift=0.33cm]bola) ++(-0.33,-0.33) arc[start angle=180, end angle=0, radius=0.33];
\draw[-] (cli.south) -- ([yshift=0.33cm]bola);
\node[rot, anchor=west, align=left] at ([xshift=0.5cm]bola) {API HTTP/JSON\\
  (\texttt{Authorization: Bearer} e \texttt{X-Justificativa};\\ pelo nginx: \texttt{/api/} $\to$ \texttt{api:8000})};
% ======================= back-end =======================
\node[comp] (api)  at (7.8,6.4) {\cmp{api}{buscador/api/ (FastAPI)}};
\draw[-] (api.north) -- ([yshift=-0.17cm]bola);
\node[comp] (agd)  at (18.7,0.6) {\cmp{agendador}{agendador.py}};
\draw[dep] ([yshift=0.15cm]api.east) -| node[rot, pos=0.25, above]{inicia a thread} ([xshift=-0.3cm]agd.north);
\draw[dep] ([xshift=0.3cm]agd.north) |- node[rot, pos=0.75, below]{importação tardia} ([yshift=-0.15cm]api.east);
% linha A
\node[comp] (ia)   at (1.9,3.0)   {\cmp{ia}{buscador/ia/}};
\node[comp] (col)  at (5.8,3.0)   {\cmp{coleta}{coleta.py}};
\node[comp] (rel)  at (9.7,3.0)   {\cmp{relatório}{buscador/relatorio/}};
\node[comp] (tip)  at (13.6,3.0)  {\cmp{tipologias}{buscador/tipologias/}};
% linha B
\node[comp] (res)  at (2.5,0.6)   {\cmp{resolução}{buscador/resolucao/}};
\node[comp] (aud)  at (11.7,0.6)  {\cmp{auditoria}{buscador/auditoria/}};
% linha C
\node[comp] (con)  at (4.1,-1.8)  {\cmp{conectores}{buscador/conectores/}};
\node[comp] (cus)  at (7.9,-1.8)  {\cmp{custódia}{buscador/custodia/}};
\node[comp] (cri)  at (11.7,-1.8) {\cmp{cripto}{cripto.py}};
\node[comp] (ace)  at (15.5,-1.8) {\cmp{acesso}{buscador/acesso/}};
% linha D (sem dependências entre módulos de domínio)
\node[comp] (gra)  at (4.1,-4.1)  {\cmp{grafo}{buscador/grafo/}};
\node[comp] (mon)  at (7.9,-4.1)  {\cmp{monitoramento}{monitoramento.py}};
\node[comp] (exp)  at (11.7,-4.1) {\cmp{exportação}{exportacao.py}};
% dependências entre módulos de domínio
\draw[dep] (ia) -- (res);
\draw[dep] ([xshift=-1.2cm]col.south) -- ([xshift=-0.3cm]res.north east);
\draw[dep] ([xshift=-0.45cm]col.south) -- ([xshift=0.6cm]con.north);
\draw[dep] ([xshift=0.6cm]col.south) -- ([xshift=-0.4cm]cus.north);
\draw[dep] (res) -- (con);
\draw[dep] (rel) -- (tip);
\draw[dep] (rel) -- (aud);
\draw[dep] ([xshift=-0.6cm]rel.south) -- ([xshift=0.6cm]cus.north);
\draw[dep] (cus) -- (cri);
\draw[dep] (ace) -- (cri);
\draw[dep] (agd) -- (aud);
\node[sub, dashed, fit={(ia)(tip)(gra)(exp)(ace) ($(ia.north)+(0,0.7)$)}] (DOM) {};
\node[titulo, font=\footnotesize\bfseries] at (DOM.north west) {Módulos de domínio (todos importados pela \texttt{api})};
\draw[dep] (api) -- node[rot, right=2pt]{usa} (api |- DOM.north);
\node[sub, fit={(api)(DOM)(agd) ($(api.north)+(0,1.0)$)}] (BE) {};
\node[titulo] at (BE.north west) {Back-end: pacote \texttt{buscador} (Python)};
% ícones de componente
\foreach \n in {telas, cmps, sess, cli, sim, api, agd, ia, col, rel, tip, res, aud, con, cus, cri, ace, gra, mon, exp}
  \pic at ([xshift=-0.36cm, yshift=-0.42cm]\n.north east) {icone};
\end{tikzpicture}}
\fonte{Elaborada pelo autor a partir das importações do código: \texttt{import} dos módulos de
\texttt{buscador/} (back-end) e dos arquivos de \texttt{src/} (front-end), além de
\texttt{nginx.conf}. Setas tracejadas: dependência (quem importa aponta para quem é importado).
Módulos utilitários (\texttt{config}, \texttt{modelos}, \texttt{db}, \texttt{seguranca},
\texttt{retencao}, \texttt{comparacao}) e importações só de funções auxiliares foram omitidos.}
\end{figure}

% Diagrama de implantação (UML 2) do buscador OSINT no Docker Compose. Elementos e fontes:
% docs/sistema/diagramas-fontes.md, seção 3.
\begin{figure}[htbp]
\centering
\caption{Diagrama de implantação do buscador OSINT (Docker Compose)}
\label{fig:uml-implantacao}
\resizebox{\textwidth}{!}{%
\begin{tikzpicture}[
  font=\small,
  no/.style={draw, fill=white, rectangle split, rectangle split parts=2, align=left,
    text width=#1, inner sep=3pt, font=\scriptsize},
  no/.default=4.9cm,
  vol/.style={draw, fill=white, align=left, text width=4.6cm, inner sep=3pt, font=\scriptsize},
  ext/.style={draw, fill=white, align=left, text width=5.7cm, inner sep=3pt, font=\scriptsize},
  cam/.style={-, thick},
  rot/.style={font=\scriptsize, fill=white, inner sep=1pt, align=center},
  moldura/.style={draw, inner sep=0.35cm},
  titulo/.style={anchor=north west, font=\footnotesize, inner sep=4pt, align=left},
  nota/.style={align=left, text width=#1, font=\scriptsize, inner sep=4pt},
  nota/.default=4.2cm
]
\newcommand{\cubo}[1]{%
  \draw (#1.north west) -- ++(0.22,0.22) -- ([xshift=0.22cm,yshift=0.22cm]#1.north east)
        -- (#1.north east);
  \draw ([xshift=0.22cm,yshift=0.22cm]#1.north east) -- ([xshift=0.22cm,yshift=0.22cm]#1.south east)
        -- (#1.south east);}
\newcommand{\dobra}[1]{%
  \draw (#1.north west) -- ([xshift=-6pt]#1.north east) -- ([yshift=-6pt]#1.north east)
        -- (#1.south east) -- (#1.south west) -- cycle;
  \draw ([xshift=-6pt]#1.north east) |- ([yshift=-6pt]#1.north east);}
% ------------------------- serviços externos (topo) -------------------------
\node[ext, anchor=north] (pncp) at (0,0.7) {%
  «serviço externo» \textbf{PNCP}\\
  \texttt{https://pncp.gov.br/api/consulta}\\ \texttt{https://pncp.gov.br/api/pncp} (termos)};
\node[ext, anchor=north] (ptr) at (6.6,0.7) {%
  «serviço externo» \textbf{Portal da Transparência}\\
  \texttt{https://api.portaldatransparencia}\\\texttt{.gov.br}; cabeçalho \texttt{chave-api-dados}};
\node[ext, dashed, anchor=north] (ia) at (13.2,0.7) {%
  «serviço» \textbf{Servidor de IA local} (opcional)\\
  Ollama em \texttt{IA\_OLLAMA\_URL}; desligado por padrão (\texttt{IA\_EXTRACAO\_ATIVA=0});
  só host local ou rede privada};
% ------------------------- host: navegador -------------------------
\node[draw, fill=white, align=center, font=\scriptsize, inner sep=3pt, anchor=north] (nav) at (0,-3.4)
  {«ambiente de execução»\\ \textbf{Navegador do usuário}};
% ------------------------- contêineres -------------------------
\node[no, anchor=north] (fe) at (0,-6.3) {%
  «contêiner» \textbf{frontend}\\ \texttt{osint-tcc-frontend}\\ imagem \texttt{nginx:alpine}
  \nodepart{second}
  usuário \texttt{nginx}; escuta \texttt{:8080}\\
  «artefato» SPA (\texttt{dist/})\\
  «artefato» \texttt{nginx.conf}: \texttt{/api/} $\to$ \texttt{api:8000}};
\node[no, anchor=north] (api) at (6.6,-6.3) {%
  «contêiner» \textbf{api}\\ \texttt{osint-tcc-api}\\ imagem \texttt{python:3.12-slim-bookworm}
  \nodepart{second}
  «artefato» \texttt{uvicorn buscador.main:app}, \texttt{:8000}; thread do agendador\\
  usuário \texttt{osint}; \texttt{cap\_drop: ALL}; \texttt{no-new-privileges}; 1 GB\\[2pt]
  \textbf{Segredos:} \texttt{AUDITORIA\_CHAVE\_HMAC}, \texttt{CUSTODIA\_CHAVE},
  \texttt{POSTGRES\_APP\_PASSWORD} (em \texttt{DATABASE\_URL}),
  \texttt{PORTAL\_TRANSPARENCIA\_CHAVE} (opcional)};
\node[nota=5.0cm, anchor=north] (ptapi) at (13.9,-6.6) {Porta da API também publicada no host:
  \texttt{127.0.0.1:8100} $\to$ \texttt{:8000}. Nenhuma porta é publicada fora de \texttt{127.0.0.1}.};
\dobra{ptapi}
\draw[dashed] (ptapi.west) -- (ptapi.west -| api.east);
\node[no, anchor=north] (pg) at (6.6,-13.9) {%
  «contêiner» \textbf{postgres}\\ \texttt{osint-tcc-postgres}\\ imagem \texttt{postgres:16-alpine}
  \nodepart{second}
  \texttt{:5432}, sem porta publicada no host; banco \texttt{osint}\\
  «artefato» \texttt{01-papel-aplicacao.sh} (montado de \texttt{./docker/postgres-init} em
  \texttt{/docker-entrypoint-initdb.d}, só leitura): cria \texttt{osint\_app}\\
  papel \texttt{osint} (dono): tabelas e gatilhos da auditoria\\
  papel \texttt{osint\_app} (aplicação): só \texttt{SELECT} e \texttt{INSERT} em
  \texttt{auditoria}, \texttt{capturas}, \texttt{alerta\_decisoes}, \texttt{anotacoes} e
  \texttt{tarefa\_eventos}; \texttt{SELECT}, \texttt{INSERT}, \texttt{UPDATE} e \texttt{DELETE}
  nas demais\\[2pt]
  \textbf{Segredos:} \texttt{POSTGRES\_PASSWORD}, \texttt{POSTGRES\_APP\_PASSWORD}};
\node[no, anchor=north] (mig) at (13.9,-14.5) {%
  «contêiner» \textbf{migracao}\\ \texttt{osint-tcc-migracao}\\ mesma imagem da \texttt{api}
  \nodepart{second}
  \texttt{python -m buscador.migrar}; roda uma vez e termina; cria tabelas e gatilhos e
  concede privilégios a \texttt{osint\_app}\\
  \texttt{cap\_drop: ALL}; \texttt{no-new-privileges}\\[2pt]
  \textbf{Segredo:} \texttt{POSTGRES\_PASSWORD}};
\node[vol, anchor=north] (vcap) at (0,-14.5) {%
  «volume» \textbf{\texttt{osint-tcc-capturas}}\\ montado em \texttt{/app/dados}:\\
  \texttt{capturas/<aa>/<sha256>.gcm} (AES-256-GCM)\\
  \texttt{ancora/auditoria.json} (âncora da trilha)};
\node[vol, anchor=north] (vpg) at (6.6,-22.9) {%
  «volume» \textbf{\texttt{osint-tcc-postgres}}\\ montado em \texttt{/var/lib/postgresql/data}};
% caminhos internos
\draw[cam] (nav) -- node[rot, right=2pt]{HTTP \texttt{127.0.0.1:8090} $\to$ \texttt{:8080}} (fe);
\coordinate (f1) at ([yshift=-0.9cm]fe.north east);
\draw[cam] (f1) -- node[rot, above]{HTTP} (f1 -| api.west);
\draw[cam] (api.south) -- node[rot, right=2pt]{SQL \texttt{:5432}, papel \texttt{osint\_app}} (pg.north);
\coordinate (m1) at ([yshift=-0.9cm]mig.north);
\draw[cam] (mig.west |- m1) -- node[rot, above]{SQL \texttt{:5432}}
  node[rot, below]{papel \texttt{osint}} (m1 -| pg.east);
\draw[cam, dashed] ([xshift=-1.6cm]api.south) -- ++(0,-0.45) -| node[rot, pos=0.25, above]{montagem} (vcap.north);
\draw[cam, dashed] (pg.south) -- node[rot, right=2pt]{montagem} (vpg.north);
% ------------------------- molduras -------------------------
\node[moldura, fit={(fe)(api)(ptapi)(mig)(pg)(vcap)(vpg) ($(fe.north)+(0,1.05)$)}] (comp) {};
\node[titulo, anchor=north east] at (comp.north east) {«ambiente de execução» \textbf{Docker Compose},
  projeto \texttt{osint-tcc}\\ rede padrão do Compose (sem rede nomeada)};
\node[moldura, fit={(comp)(nav) ($(nav.north)+(0,0.75)$)}] (host) {};
\node[titulo] at (host.north west) {«dispositivo» \textbf{Máquina host} (Docker Engine)};
\foreach \n in {fe, api, mig, pg, host, pncp, ptr, ia} \cubo{\n}
% api -> serviços externos (passam pela fronteira do host)
\coordinate (bus) at (0,-1.75);
\draw[cam] (api.north) -- (api.north |- bus) -| node[rot, pos=0.6, right]{HTTPS} (pncp.south);
\draw[cam] (api.north |- bus) -- node[rot, right, pos=0.4]{HTTPS} (ptr.south);
\draw[cam, dashed] (api.north |- bus) -| node[rot, pos=0.4, above]{HTTP} (ia.south);
% ------------------------- notas -------------------------
\node[nota=5.7cm, anchor=north] (n55) at (13.9,-22.4) {Só com o arquivo extra
  \texttt{docker-compose.postgres-porta.yml}: \texttt{5432} publicada em \texttt{127.0.0.1:55432}.};
\dobra{n55}
\coordinate (q55) at ([xshift=0.5cm,yshift=0.7cm]pg.south east);
\draw[dashed] (n55.west) -| (q55) -- ([yshift=0.7cm]pg.south east);
\node[nota=4.3cm, anchor=north] (nord) at (0,-22.4) {Ordem de subida (\texttt{depends\_on}):
  postgres saudável $\to$ migracao concluída $\to$ api saudável $\to$ frontend.};
\dobra{nord}
\end{tikzpicture}}
\fonte{Elaborada pelo autor a partir de \texttt{docker-compose.yml},
\texttt{docker-compose.postgres-porta.yml}, \texttt{Dockerfile} e
\texttt{docker/postgres-init/01-papel-aplicacao.sh} do back-end, de \texttt{Dockerfile} e
\texttt{nginx.conf} do front-end e de \texttt{buscador/db.py} (privilégios do papel da aplicação).
Segredos: só os nomes das variáveis, definidas no arquivo \texttt{.env}.}
\end{figure}


\textbf{Fontes de dados.} A coleta é só passiva, por conectores com nome e
versão: dados abertos do CNPJ (arquivos no layout da Receita Federal), API de
consulta e termos de contrato do PNCP, API do Portal da Transparência (CEIS,
CNEP, CEPIM, CEAF, servidores e convênios, com chave do usuário), dados abertos do TSE (prestação de
contas e candidatos) e editais em texto ou em PDF com camada de texto. Os
fatos usados pelas regras (limites
da Lei 14.133/2021, decreto de atualização de 2026, campos das APIs e códigos
de situação cadastral) foram conferidos nas fontes oficiais
(\texttt{docs/sistema/fontes-horizonte2.md}).

\textbf{Editais em PDF.} O conector de editais lê também PDF com camada de
texto, com a biblioteca \texttt{pypdf} (versão 7.1.0, licença BSD-3-Clause).
A custódia é sobre os bytes originais: a captura guarda o PDF como veio, e o
SHA-256 registrado é o do arquivo. O texto é derivado. A rota de texto, a
tipologia T10 e a extração por IA extraem o texto de novo a partir da
captura, e as posições de um trecho referem-se ao texto extraído, e não ao
arquivo. Para que essa referência seja reproduzível, a captura registra o
número de páginas, o extrator (biblioteca e versão) e o SHA-256 do texto
extraído na coleta. Se outra versão da biblioteca extrair outro texto, a
diferença é sinalizada. A extração roda num processo separado, sem as
variáveis de ambiente da API, e é encerrada no prazo. Há limites de tamanho
(50~MiB), de páginas (300), de caracteres (5 milhões) e de tempo (60~s). O
código só extrai o texto das páginas: não executa JavaScript nem ações do
documento e não abre anexos. PDF criptografado, corrompido ou acima dos
limites é recusado, com o motivo. PDF escaneado fica na custódia, marcado como
``sem camada de texto'', porque não há OCR. O conector não baixa o edital do
PNCP: o investigador o põe na pasta de ingestão.

\textbf{Recorte do CNPJ.} O conector do CNPJ copia cada arquivo para a
custódia e transforma cada linha em entidade do caso, com tetos de 16~MiB por
arquivo e de 100 mil registros por coleta. A base completa de 2026-11 tem 37
arquivos ZIP, com 7,87~GB compactados e, por proporção, mais de 26~GB
descompactados nas três tabelas lidas. Por isso ela não passa pelo conector.
Um \textit{script} (\texttt{scripts/recortar\_cnpj.py}) lê os ZIPs oficiais
em fluxo, sem descompactar em disco, e grava só as linhas dos CNPJs do caso,
com um manifesto que registra o SHA-256 de cada ZIP lido. O conector lê esse
recorte normalmente. O recorte é dado derivado: a captura registra o SHA-256
do recorte, e o manifesto o liga aos ZIPs de origem. Na medição, o
\textit{script} leu cerca de 110~MB/s.

% ===========================================================================
\section{Casos de uso}
\label{sec:desenv-casos-uso}

A Figura~\ref{fig:uml-casos-de-uso} liga os atores às funções do sistema,
agrupadas pelos épicos do \textit{backlog}. Os papéis vêm do código: perfis
globais (analista, delegado e administrador) e papéis no caso (responsável,
investigador e leitor). O código não tem perfis separados de auditor e de
perito: a trilha de auditoria é listada pelos perfis delegado e
administrador, e a recoleta exige o papel de investigador no caso.

% Diagrama de casos de uso (UML 2) do buscador OSINT. Elementos e fontes no código:
% docs/sistema/diagramas-fontes.md, seção 1.
\begin{figure}[p]
\centering
\caption{Diagrama de casos de uso do buscador OSINT, agrupados pelos épicos do backlog}
\label{fig:uml-casos-de-uso}
\resizebox{0.95\textwidth}{!}{%
\begin{tikzpicture}[
  font=\small,
  uc/.style={ellipse, draw, text width=4.3cm, align=center, inner sep=-1pt,
    minimum height=1.05cm, font=\small,
    execute at begin node={\hyphenpenalty=10000\exhyphenpenalty=10000}},
  pacote/.style={draw, inner sep=0.22cm},
  aba/.style={draw, anchor=south west, font=\scriptsize\bfseries, inner sep=2pt},
  atorcx/.style={minimum width=0.7cm, minimum height=1.3cm, inner sep=0pt},
  rotulo/.style={align=center, font=\small, inner sep=1pt},
  assoc/.style={-},
  gen/.style={-{Triangle[open, length=7pt, width=7pt]}},
  nota/.style={align=left, text width=5.2cm, font=\scriptsize, inner sep=4pt},
  ator/.pic={
    \draw (0,0.45) circle (0.15);
    \draw (0,0.3) -- (0,-0.15) (-0.28,0.15) -- (0.28,0.15)
          (0,-0.15) -- (-0.22,-0.55) (0,-0.15) -- (0.22,-0.55);}
]
% ---------------- casos de uso ----------------
\def\xe{5.3}
\def\xd{12.4}
% E2
\node[uc] (gra) at (\xe,-0.6)   {Ver grafo de vínculos e caminho entre entidades};
% E4
\node[uc] (ver) at (\xe,-2.85)  {Verificar integridade da captura};
\node[uc] (rel) at (\xe,-4.3)   {Gerar relatório e exportar o grafo};
\node[uc] (rec) at (\xe,-5.75)  {Recoletar e comparar capturas};
% E3
\node[uc] (val) at (\xe,-8.0)   {Ver alertas por severidade};
\node[uc] (exe) at (\xe,-9.45)  {Executar tipologias};
\node[uc] (dec) at (\xe,-10.9)  {Confirmar ou descartar alerta, com justificativa};
\node[uc] (ia)  at (\xe,-12.35) {Extrair com IA local e decidir sugestão};
% E1
\node[uc] (col) at (\xe,-14.6)  {Coletar de fonte aberta};
% E7
\node[uc] (mon) at (\xe,-16.85) {Marcar entidade para monitoramento};
\node[uc] (exm) at (\xe,-18.3)  {Executar monitoramento e gerar avisos};
% E5
\node[uc] (bus) at (\xe,-20.55) {Buscar entidades por nome, CNPJ ou CPF};
\node[uc] (ano) at (\xe,-22.0)  {Anotar e dividir tarefas};
\node[uc] (cri) at (\xd,-21.3)  {Criar caso com finalidade e prazo de retenção};
% E6
\node[uc] (mem) at (\xe,-26.4)  {Gerenciar membros e papéis do caso};
\node[uc] (ent) at (\xd,-24.25) {Entrar com senha e código TOTP};
\node[uc] (cad) at (\xd,-25.7)  {Verificar cadeia da auditoria};
\node[uc] (usu) at (\xd,-27.15) {Gerenciar usuários e perfis};
\node[uc] (tri) at (\xd,-28.6)  {Consultar trilha de auditoria};
% ---------------- pacotes (épicos) ----------------
\node[pacote, fit=(gra)] (E2) {};
\node[aba] at (E2.north west) {E2. Entidades e vínculos};
\node[pacote, fit=(ver)(rec)] (E4) {};
\node[aba] at (E4.north west) {E4. Prova e cadeia de custódia};
\node[pacote, fit=(val)(ia)] (E3) {};
\node[aba] at (E3.north west) {E3. Alertas e análise assistida};
\node[pacote, fit=(col)] (E1) {};
\node[aba] at (E1.north west) {E1. Coleta de dados abertos};
\node[pacote, fit=(mon)(exm)] (E7) {};
\node[aba] at (E7.north west) {E7. Monitoramento e evolução};
\node[pacote, fit=(bus)(ano)(cri)] (E5) {};
\node[aba] at (E5.north west) {E5. Casos e busca};
\node[pacote, fit=(mem)(ent)(tri)] (E6) {};
\node[aba] at (E6.north west) {E6. Segurança, auditoria e conformidade};
% ---------------- fronteira do sistema ----------------
\node[draw, thick, inner xsep=0.35cm, inner ysep=0.3cm,
      fit={(E2.north west) ($(E2.north west)+(0,0.75)$) (E6.south west) (E6.east)}] (sis) {};
\node[anchor=north, font=\small\bfseries] at (sis.north) {Buscador OSINT};
% ---------------- atores ----------------
\node[atorcx] (invest) at (0.6,-11.0) {}; \pic at (invest) {ator};
\node[rotulo, left=1pt of invest] (investR) {Investigador\\{\scriptsize(papel no caso)}};
\node[atorcx] (resp) at (0.6,-26.4) {}; \pic at (resp) {ator};
\node[rotulo, left=1pt of resp] (respR) {Delegado\\responsável\\{\scriptsize(papel \texttt{responsavel})}};
\node[atorcx] (leitor) at (17.1,-4.0) {}; \pic at (leitor) {ator};
\node[rotulo, right=1pt of leitor] (leitorR) {Leitor\\{\scriptsize(papel no caso)}};
\node[atorcx] (agend) at (17.1,-18.3) {}; \pic at (agend) {ator};
\node[rotulo, right=1pt of agend] (agendR) {Agendador\\{\scriptsize(ator sistema)}};
\node[atorcx] (usuario) at (17.1,-22.6) {}; \pic at (usuario) {ator};
\node[rotulo, right=1pt of usuario] (usuarioR) {Usuário\\autenticado\\{\scriptsize(qualquer perfil)}};
\node[atorcx] (admin) at (17.1,-26.0) {}; \pic at (admin) {ator};
\node[rotulo, right=1pt of admin] (adminR) {Administrador\\{\scriptsize(perfil \texttt{admin})}};
\node[atorcx] (deleg) at (21.4,-25.6) {}; \pic at (deleg) {ator};
\node[rotulo, below=0pt of deleg] (delegR) {Delegado\\{\scriptsize(perfil)}};
% ---------------- generalizações ----------------
\draw[gen] (resp.north) -- (invest.south);
\draw[gen] (admin.north) -- (usuario.south);
\draw (deleg.north) |- (17.1,-24.7);
% Investigador especializa Leitor: a linha contorna a fronteira por cima.
\draw[gen] (invest.north) -- (0.6,1.5) -| (leitor.north);
% ---------------- associações ----------------
\foreach \u in {gra, ver, rel, val, bus} \draw[assoc] (leitor.west) -- (\u.east);
\foreach \u in {rec, exe, dec, ia, col, mon, exm, ano} \draw[assoc] (invest.east) -- (\u.west);
\draw[assoc] (resp.east) -- (mem.west);
\draw[assoc] (agend.west) -- (exm.east);
\foreach \u in {cri, ent, cad} \draw[assoc] (usuario.west) -- (\u.east);
\draw[assoc] (admin.west) -- (usu.east);
\draw[assoc] (admin.west) -- (tri.east);
\draw[assoc] (delegR.south) |- (tri.east);
% ---------------- nota ----------------
\node[nota, text width=17.5cm, anchor=north west] (nt) at ($(sis.south west)+(-1.6,-0.5)$) {%
  \textbf{Perfis globais} (\texttt{usuarios.perfil}): analista, delegado e admin; Delegado e
  Administrador especializam Usuário autenticado. \textbf{Papéis no caso}
  (\texttt{caso\_membros}): leitor $<$ investigador $<$ responsável; quem cria o caso é o primeiro
  responsável, com qualquer perfil. Não há perfil de auditor nem de perito: a trilha é consultada
  pelos perfis delegado e admin, e a recoleta (H14) exige o papel de investigador. Toda ação sobre
  dado de caso exige justificativa (\texttt{X-Justificativa}) e acrescenta uma entrada à trilha de
  auditoria.};
\draw (nt.north west) -- ([xshift=-7pt]nt.north east) -- ([yshift=-7pt]nt.north east)
      -- (nt.south east) -- (nt.south west) -- cycle;
\draw ([xshift=-7pt]nt.north east) |- ([yshift=-7pt]nt.north east);
\end{tikzpicture}}
\fonte{Elaborada pelo autor a partir do código do back-end (\texttt{buscador/api/deps.py},
\texttt{buscador/api/rotas\_*.py}, \texttt{buscador/agendador.py}, \texttt{buscador/config.py}) e do
backlog (\texttt{docs/backlog/backlog.md}). Só aparecem casos de uso com rota implementada.}
\end{figure}


% ===========================================================================
\section{Modelo de dados}
\label{sec:desenv-dados}

O banco do dado aberto tem 21 tabelas, e o módulo de importação sob ordem
judicial (RF10) tem outras 4, numa base separada. As
Figuras~\ref{fig:modelo-coleta}, \ref{fig:modelo-analise},
\ref{fig:modelo-acesso} e~\ref{fig:modelo-judicial} mostram o modelo em
quatro áreas, com as multiplicidades de cada relação. A
Tabela~\ref{tab:relacoes} lista as 42 chaves estrangeiras e as 12 referências
sem chave estrangeira, com
a cardinalidade de cada uma. O Apêndice~\ref{ap:dicionario} traz o dicionário
de dados, com a descrição de cada coluna. As figuras e as tabelas são geradas
a partir do modelo do próprio código, para não divergirem do sistema.

O caso é o centro do modelo: quase toda tabela tem a coluna
\texttt{caso\_id}, numa relação 1:N com \texttt{casos}. Quatro tabelas são
só de acréscimo, protegidas por gatilhos do banco contra alteração e
exclusão: \texttt{auditoria}, \texttt{alerta\_decisoes}, \texttt{anotacoes} e
\texttt{tarefa\_eventos}. Nas capturas, o papel do banco usado pela aplicação
só pode ler e inserir.

% Gerado por docs/sistema/gera_modelo_tex.py a partir de docs/sistema/modelo.json. Não editar à mão.
\begin{figure}[p]
\centering
\caption{Modelo de dados: Coleta, custódia e grafo}
\label{fig:modelo-coleta}
\resizebox{0.95\textwidth}{!}{%
\begin{tikzpicture}[
tabela/.style={rectangle split, rectangle split parts=2, draw, rounded corners=1pt, font=\footnotesize, align=left, inner sep=3pt, anchor=north, fill=white},
fk/.style={-, thick}, logica/.style={-, dashed}, mult/.style={font=\scriptsize, fill=white, inner sep=1pt}]
\node[tabela] (coletas) at (-6.0,0) {\textbf{coletas}\nodepart{second}\begin{tabular}{@{}l@{\hspace{4pt}}l@{}}{\scriptsize PK }\underline{id} & {\scriptsize texto(32)} \\ {\scriptsize FK }caso\_id & {\scriptsize texto(32)} \\ {\scriptsize }fonte & {\scriptsize texto(50)} \\ {\scriptsize }parametros & {\scriptsize JSON} \\ {\scriptsize }conector\_nome & {\scriptsize texto(50)} \\ {\scriptsize }conector\_versao & {\scriptsize texto(20)} \\ {\scriptsize }situacao & {\scriptsize texto(15)} \\ {\scriptsize }tentativas & {\scriptsize inteiro} \\ {\scriptsize }motivo & {\scriptsize texto} \\ {\scriptsize }registros & {\scriptsize inteiro} \\ {\scriptsize }avisos & {\scriptsize JSON} \\ {\scriptsize }usuario & {\scriptsize texto(100)} \\ {\scriptsize }iniciado\_em & {\scriptsize texto(40)} \\ {\scriptsize }concluido\_em & {\scriptsize texto(40)?} \\ {\scriptsize FK }recoleta\_de & {\scriptsize texto(32)?}\end{tabular}};
\node[tabela] (capturas) at (-6.0,-8.4) {\textbf{capturas}\nodepart{second}\begin{tabular}{@{}l@{\hspace{4pt}}l@{}}{\scriptsize PK }\underline{id} & {\scriptsize texto(32)} \\ {\scriptsize FK }caso\_id & {\scriptsize texto(32)} \\ {\scriptsize }fonte & {\scriptsize texto(50)} \\ {\scriptsize }url\_ou\_arquivo & {\scriptsize texto} \\ {\scriptsize }coletado\_em & {\scriptsize texto(40)} \\ {\scriptsize }sha256 & {\scriptsize texto(64)} \\ {\scriptsize }tamanho\_bytes & {\scriptsize inteiro} \\ {\scriptsize }caminho\_relativo & {\scriptsize texto(200)} \\ {\scriptsize }metadados & {\scriptsize JSON} \\ {\scriptsize }conector\_nome & {\scriptsize texto(50)} \\ {\scriptsize }conector\_versao & {\scriptsize texto(20)} \\ {\scriptsize FK }coleta\_id & {\scriptsize texto(32)?}\end{tabular}};
\node[tabela] (entidades) at (1.0,-8.4) {\textbf{entidades}\nodepart{second}\begin{tabular}{@{}l@{\hspace{4pt}}l@{}}{\scriptsize PK }\underline{id} & {\scriptsize texto(32)} \\ {\scriptsize FK }caso\_id & {\scriptsize texto(32)} \\ {\scriptsize }tipo & {\scriptsize texto(20)} \\ {\scriptsize }chave & {\scriptsize texto(300)} \\ {\scriptsize }rotulo & {\scriptsize texto(300)} \\ {\scriptsize }identificador & {\scriptsize texto(100)} \\ {\scriptsize }atributos & {\scriptsize JSON} \\ {\scriptsize FK }captura\_id & {\scriptsize texto(32)}\end{tabular}};
\node[tabela] (entidadefontes) at (-2.6,-16.0) {\textbf{entidade\_fontes}\nodepart{second}\begin{tabular}{@{}l@{\hspace{4pt}}l@{}}{\scriptsize PK }\underline{id} & {\scriptsize inteiro} \\ {\scriptsize FK }entidade\_id & {\scriptsize texto(32)} \\ {\scriptsize FK }captura\_id & {\scriptsize texto(32)} \\ {\scriptsize }papel & {\scriptsize texto(50)}\end{tabular}};
\node[tabela] (arestas) at (8.0,-15.0) {\textbf{arestas}\nodepart{second}\begin{tabular}{@{}l@{\hspace{4pt}}l@{}}{\scriptsize PK }\underline{id} & {\scriptsize texto(32)} \\ {\scriptsize FK }caso\_id & {\scriptsize texto(32)} \\ {\scriptsize FK }origem\_id & {\scriptsize texto(32)} \\ {\scriptsize FK }destino\_id & {\scriptsize texto(32)} \\ {\scriptsize }tipo & {\scriptsize texto(30)} \\ {\scriptsize }chave & {\scriptsize texto(200)} \\ {\scriptsize }atributos & {\scriptsize JSON} \\ {\scriptsize FK }captura\_id & {\scriptsize texto(32)}\end{tabular}};
\node[tabela] (sugestoes) at (8.0,-4.0) {\textbf{sugestoes}\nodepart{second}\begin{tabular}{@{}l@{\hspace{4pt}}l@{}}{\scriptsize PK }\underline{id} & {\scriptsize texto(32)} \\ {\scriptsize FK }caso\_id & {\scriptsize texto(32)} \\ {\scriptsize FK }entidade\_id & {\scriptsize texto(32)} \\ {\scriptsize FK }candidata\_id & {\scriptsize texto(32)} \\ {\scriptsize }score & {\scriptsize real} \\ {\scriptsize }motivo & {\scriptsize texto} \\ {\scriptsize }criado\_em & {\scriptsize texto(40)}\end{tabular}};
\draw[fk] (coletas.south west) to[out=180,in=180,looseness=1.6] node[mult,pos=0.12]{0..*} node[mult,pos=0.88]{0..1} node[mult,pos=0.5,left]{recoleta\_de} (coletas.north west);
\draw[fk] (capturas) to[] node[mult,pos=0.2]{0..*} node[mult,pos=0.8]{0..1} (coletas);
\draw[fk] (entidades) to[] node[mult,pos=0.2]{0..*} node[mult,pos=0.8]{1} (capturas);
\draw[fk] (arestas) to[] node[mult,pos=0.2]{0..*} node[mult,pos=0.8]{1} (entidades);
\draw[fk] (arestas) to[bend left=18] node[mult,pos=0.2]{0..*} node[mult,pos=0.5]{destino\_id} node[mult,pos=0.8]{1} (entidades);
\draw[fk] (arestas) to[] node[mult,pos=0.2]{0..*} node[mult,pos=0.8]{1} (capturas);
\draw[fk] (entidadefontes) to[] node[mult,pos=0.2]{0..*} node[mult,pos=0.8]{1} (entidades);
\draw[fk] (entidadefontes) to[] node[mult,pos=0.2]{0..*} node[mult,pos=0.8]{1} (capturas);
\draw[fk] (sugestoes) to[] node[mult,pos=0.2]{0..*} node[mult,pos=0.8]{1} (entidades);
\draw[fk] (sugestoes) to[bend left=18] node[mult,pos=0.2]{0..*} node[mult,pos=0.5]{candidata\_id} node[mult,pos=0.8]{1} (entidades);
\end{tikzpicture}}
\fonte{Gerada pelo autor a partir do modelo SQLAlchemy do back-end (\texttt{buscador/modelos.py}). PK: chave primária; FK: chave estrangeira; ``?'': aceita nulo. Linha cheia: chave estrangeira; linha tracejada: referência sem chave estrangeira. Multiplicidade junto de cada ponta. Tabelas mostradas só com as chaves: colunas completas no Apêndice~\ref{ap:dicionario}. Para legibilidade, a ligação de cada tabela com \texttt{casos} pela coluna \texttt{caso\_id} (1:N) não aparece; ela está na Tabela~\ref{tab:relacoes}.}
\end{figure}

\begin{figure}[p]
\centering
\caption{Modelo de dados: Alertas, monitoramento e extração por IA}
\label{fig:modelo-analise}
\resizebox{0.95\textwidth}{!}{%
\begin{tikzpicture}[
tabela/.style={rectangle split, rectangle split parts=2, draw, rounded corners=1pt, font=\footnotesize, align=left, inner sep=3pt, anchor=north, fill=white},
fk/.style={-, thick}, logica/.style={-, dashed}, mult/.style={font=\scriptsize, fill=white, inner sep=1pt}]
\node[tabela] (alertas) at (-7.0,0) {\textbf{alertas}\nodepart{second}\begin{tabular}{@{}l@{\hspace{4pt}}l@{}}{\scriptsize PK }\underline{id} & {\scriptsize texto(32)} \\ {\scriptsize FK }caso\_id & {\scriptsize texto(32)} \\ {\scriptsize }tipologia & {\scriptsize texto(5)} \\ {\scriptsize }titulo & {\scriptsize texto(300)} \\ {\scriptsize }motivo & {\scriptsize texto} \\ {\scriptsize }severidade & {\scriptsize texto(10)} \\ {\scriptsize }entidades & {\scriptsize JSON} \\ {\scriptsize }evidencias & {\scriptsize JSON} \\ {\scriptsize }parametros & {\scriptsize JSON} \\ {\scriptsize }criado\_em & {\scriptsize texto(40)} \\ {\scriptsize }chave & {\scriptsize texto(64)} \\ {\scriptsize }ativo & {\scriptsize lógico} \\ {\scriptsize }situacao & {\scriptsize texto(12)} \\ {\scriptsize }decidido\_por & {\scriptsize texto(100)?} \\ {\scriptsize }decidido\_em & {\scriptsize texto(40)?} \\ {\scriptsize }justificativa\_decisao & {\scriptsize texto?}\end{tabular}};
\node[tabela] (alertadecisoes) at (-7.0,-9.6) {\textbf{alerta\_decisoes}\nodepart{second}\begin{tabular}{@{}l@{\hspace{4pt}}l@{}}{\scriptsize PK }\underline{id} & {\scriptsize inteiro} \\ {\scriptsize FK }alerta\_id & {\scriptsize texto(32)} \\ {\scriptsize FK }caso\_id & {\scriptsize texto(32)} \\ {\scriptsize }decisao & {\scriptsize texto(12)} \\ {\scriptsize }justificativa & {\scriptsize texto} \\ {\scriptsize }usuario & {\scriptsize texto(100)} \\ {\scriptsize }perfil & {\scriptsize texto(20)} \\ {\scriptsize }decidido\_em & {\scriptsize texto(40)}\end{tabular}};
\node[tabela] (avisos) at (0.0,-0.6) {\textbf{avisos}\nodepart{second}\begin{tabular}{@{}l@{\hspace{4pt}}l@{}}{\scriptsize PK }\underline{id} & {\scriptsize texto(32)} \\ {\scriptsize FK }caso\_id & {\scriptsize texto(32)} \\ {\scriptsize FK }monitoramento\_id & {\scriptsize texto(32)} \\ {\scriptsize FK }execucao\_id & {\scriptsize texto(32)} \\ {\scriptsize FK }entidade\_id & {\scriptsize texto(32)} \\ {\scriptsize }tipo & {\scriptsize texto(30)} \\ {\scriptsize }titulo & {\scriptsize texto(300)} \\ {\scriptsize }descricao & {\scriptsize texto} \\ {\scriptsize }detalhes & {\scriptsize JSON} \\ {\scriptsize }captura\_ids & {\scriptsize JSON} \\ {\scriptsize }criado\_em & {\scriptsize texto(40)} \\ {\scriptsize }lido\_em & {\scriptsize texto(40)?} \\ {\scriptsize }lido\_por & {\scriptsize texto(100)?}\end{tabular}};
\node[tabela] (monitoramentos) at (7.0,0) {\textbf{monitoramentos}\nodepart{second}\begin{tabular}{@{}l@{\hspace{4pt}}l@{}}{\scriptsize PK }\underline{id} & {\scriptsize texto(32)} \\ {\scriptsize FK }caso\_id & {\scriptsize texto(32)} \\ {\scriptsize FK }entidade\_id & {\scriptsize texto(32)} \\ {\scriptsize }periodicidade\_horas & {\scriptsize inteiro} \\ {\scriptsize }ativo & {\scriptsize lógico} \\ {\scriptsize }motivo\_inativo & {\scriptsize texto} \\ {\scriptsize }justificativa & {\scriptsize texto} \\ {\scriptsize }criado\_por & {\scriptsize texto(100)} \\ {\scriptsize }criado\_em & {\scriptsize texto(40)} \\ {\scriptsize }proxima\_execucao\_em & {\scriptsize texto(40)} \\ {\scriptsize }ultima\_execucao\_em & {\scriptsize texto(40)?}\end{tabular}};
\node[tabela] (monitoramentoexecucoes) at (7.0,-7.8) {\textbf{monitoramento\_execucoes}\nodepart{second}\begin{tabular}{@{}l@{\hspace{4pt}}l@{}}{\scriptsize PK }\underline{id} & {\scriptsize texto(32)} \\ {\scriptsize FK }monitoramento\_id & {\scriptsize texto(32)} \\ {\scriptsize FK }caso\_id & {\scriptsize texto(32)} \\ {\scriptsize }origem & {\scriptsize texto(12)} \\ {\scriptsize }usuario & {\scriptsize texto(100)} \\ {\scriptsize }iniciado\_em & {\scriptsize texto(40)} \\ {\scriptsize }concluido\_em & {\scriptsize texto(40)?} \\ {\scriptsize }situacao & {\scriptsize texto(20)} \\ {\scriptsize }coletas & {\scriptsize JSON} \\ {\scriptsize }avisos\_gerados & {\scriptsize inteiro} \\ {\scriptsize }motivo & {\scriptsize texto}\end{tabular}};
\node[tabela] (entidades) at (0.0,-10.4) {\textbf{entidades}\nodepart{second}\begin{tabular}{@{}l@{\hspace{4pt}}l@{}}{\scriptsize PK }\underline{id} & {\scriptsize texto(32)} \\ {\scriptsize FK }caso\_id & {\scriptsize texto(32)} \\ {\scriptsize FK }captura\_id & {\scriptsize texto(32)}\end{tabular}};
\node[tabela] (capturas) at (-7.0,-15.0) {\textbf{capturas}\nodepart{second}\begin{tabular}{@{}l@{\hspace{4pt}}l@{}}{\scriptsize PK }\underline{id} & {\scriptsize texto(32)} \\ {\scriptsize FK }caso\_id & {\scriptsize texto(32)} \\ {\scriptsize FK }coleta\_id & {\scriptsize texto(32)?}\end{tabular}};
\node[tabela] (extracoesia) at (0.0,-14.6) {\textbf{extracoes\_ia}\nodepart{second}\begin{tabular}{@{}l@{\hspace{4pt}}l@{}}{\scriptsize PK }\underline{id} & {\scriptsize texto(32)} \\ {\scriptsize FK }caso\_id & {\scriptsize texto(32)} \\ {\scriptsize FK }captura\_id & {\scriptsize texto(32)}\end{tabular}};
\node[tabela] (sugestoesia) at (7.0,-14.6) {\textbf{sugestoes\_ia}\nodepart{second}\begin{tabular}{@{}l@{\hspace{4pt}}l@{}}{\scriptsize PK }\underline{id} & {\scriptsize texto(32)} \\ {\scriptsize FK }caso\_id & {\scriptsize texto(32)} \\ {\scriptsize FK }extracao\_id & {\scriptsize texto(32)} \\ {\scriptsize FK }captura\_id & {\scriptsize texto(32)} \\ {\scriptsize }tipo & {\scriptsize texto(20)} \\ {\scriptsize }valor & {\scriptsize texto(500)} \\ {\scriptsize }campos & {\scriptsize JSON} \\ {\scriptsize }trecho & {\scriptsize texto} \\ {\scriptsize }posicao\_inicio & {\scriptsize inteiro} \\ {\scriptsize }posicao\_fim & {\scriptsize inteiro} \\ {\scriptsize }confianca & {\scriptsize real?} \\ {\scriptsize }situacao & {\scriptsize texto(10)} \\ {\scriptsize }decidido\_por & {\scriptsize texto(100)?} \\ {\scriptsize }decidido\_em & {\scriptsize texto(40)?} \\ {\scriptsize }justificativa\_decisao & {\scriptsize texto?} \\ {\scriptsize FK }entidade\_id & {\scriptsize texto(32)?}\end{tabular}};
\draw[fk] (alertadecisoes) to[] node[mult,pos=0.2]{0..*} node[mult,pos=0.8]{1} (alertas);
\draw[fk] (entidades) to[] node[mult,pos=0.2]{0..*} node[mult,pos=0.8]{1} (capturas);
\draw[fk] (extracoesia) to[] node[mult,pos=0.2]{0..*} node[mult,pos=0.8]{1} (capturas);
\draw[fk] (monitoramentos) to[] node[mult,pos=0.2]{0..*} node[mult,pos=0.8]{1} (entidades);
\draw[fk] (sugestoesia) to[] node[mult,pos=0.2]{0..*} node[mult,pos=0.8]{1} (extracoesia);
\draw[fk] (sugestoesia) to[] node[mult,pos=0.2]{0..*} node[mult,pos=0.8]{1} (capturas);
\draw[fk] (sugestoesia) to[] node[mult,pos=0.2]{0..*} node[mult,pos=0.8]{0..1} (entidades);
\draw[fk] (monitoramentoexecucoes) to[] node[mult,pos=0.2]{0..*} node[mult,pos=0.8]{1} (monitoramentos);
\draw[fk] (avisos) to[] node[mult,pos=0.2]{0..*} node[mult,pos=0.8]{1} (monitoramentos);
\draw[fk] (avisos) to[] node[mult,pos=0.2]{0..*} node[mult,pos=0.8]{1} (monitoramentoexecucoes);
\draw[fk] (avisos) to[] node[mult,pos=0.2]{0..*} node[mult,pos=0.8]{1} (entidades);
\draw[logica] (alertas) -- node[mult,pos=0.5]{entidades} (entidades);
\draw[logica] (alertas) -- node[mult,pos=0.5]{evidencias} (capturas);
\draw[logica] (avisos) -- node[mult,pos=0.5]{captura\_ids} (capturas);
\end{tikzpicture}}
\fonte{Gerada pelo autor a partir do modelo SQLAlchemy do back-end (\texttt{buscador/modelos.py}). PK: chave primária; FK: chave estrangeira; ``?'': aceita nulo. Linha cheia: chave estrangeira; linha tracejada: referência sem chave estrangeira. Multiplicidade junto de cada ponta. Tabelas mostradas só com as chaves: colunas completas no Apêndice~\ref{ap:dicionario}. Para legibilidade, a ligação de cada tabela com \texttt{casos} pela coluna \texttt{caso\_id} (1:N) não aparece; ela está na Tabela~\ref{tab:relacoes}.}
\end{figure}

\begin{figure}[p]
\centering
\caption{Modelo de dados: Acesso, equipe e auditoria}
\label{fig:modelo-acesso}
\resizebox{0.95\textwidth}{!}{%
\begin{tikzpicture}[
tabela/.style={rectangle split, rectangle split parts=2, draw, rounded corners=1pt, font=\footnotesize, align=left, inner sep=3pt, anchor=north, fill=white},
fk/.style={-, thick}, logica/.style={-, dashed}, mult/.style={font=\scriptsize, fill=white, inner sep=1pt}]
\node[tabela] (casos) at (0,0) {\textbf{casos}\nodepart{second}\begin{tabular}{@{}l@{\hspace{4pt}}l@{}}{\scriptsize PK }\underline{id} & {\scriptsize texto(32)} \\ {\scriptsize }titulo & {\scriptsize texto(300)} \\ {\scriptsize }descricao & {\scriptsize texto} \\ {\scriptsize }numero\_procedimento & {\scriptsize texto(100)} \\ {\scriptsize }criado\_por & {\scriptsize texto(100)} \\ {\scriptsize }criado\_em & {\scriptsize texto(40)} \\ {\scriptsize }finalidade & {\scriptsize texto} \\ {\scriptsize }retencao\_ate & {\scriptsize texto(10)?}\end{tabular}};
\node[tabela] (casomembros) at (-6.5,0) {\textbf{caso\_membros}\nodepart{second}\begin{tabular}{@{}l@{\hspace{4pt}}l@{}}{\scriptsize PK }\underline{id} & {\scriptsize inteiro} \\ {\scriptsize FK }caso\_id & {\scriptsize texto(32)} \\ {\scriptsize }usuario & {\scriptsize texto(100)} \\ {\scriptsize }papel & {\scriptsize texto(15)} \\ {\scriptsize }adicionado\_por & {\scriptsize texto(100)} \\ {\scriptsize }adicionado\_em & {\scriptsize texto(40)}\end{tabular}};
\node[tabela] (usuarios) at (-6.5,-7.0) {\textbf{usuarios}\nodepart{second}\begin{tabular}{@{}l@{\hspace{4pt}}l@{}}{\scriptsize PK }\underline{id} & {\scriptsize texto(32)} \\ {\scriptsize }usuario & {\scriptsize texto(100)} \\ {\scriptsize }nome & {\scriptsize texto(200)} \\ {\scriptsize }perfil & {\scriptsize texto(20)} \\ {\scriptsize }senha\_hash & {\scriptsize texto(300)} \\ {\scriptsize }ativo & {\scriptsize lógico} \\ {\scriptsize }mfa\_segredo\_cifrado & {\scriptsize texto?} \\ {\scriptsize }mfa\_ativo & {\scriptsize lógico} \\ {\scriptsize }mfa\_confirmado\_em & {\scriptsize texto(40)?} \\ {\scriptsize }mfa\_ultimo\_passo & {\scriptsize inteiro?} \\ {\scriptsize }falhas\_seguidas & {\scriptsize inteiro} \\ {\scriptsize }bloqueado\_ate & {\scriptsize texto(40)?} \\ {\scriptsize }criado\_por & {\scriptsize texto(100)} \\ {\scriptsize }criado\_em & {\scriptsize texto(40)} \\ {\scriptsize }senha\_alterada\_em & {\scriptsize texto(40)}\end{tabular}};
\node[tabela] (sessoes) at (-6.5,-15.4) {\textbf{sessoes}\nodepart{second}\begin{tabular}{@{}l@{\hspace{4pt}}l@{}}{\scriptsize PK }\underline{id} & {\scriptsize texto(32)} \\ {\scriptsize }token\_resumo & {\scriptsize texto(64)} \\ {\scriptsize FK }conta\_id & {\scriptsize texto(32)} \\ {\scriptsize }etapa & {\scriptsize texto(15)} \\ {\scriptsize }criada\_em & {\scriptsize texto(40)} \\ {\scriptsize }expira\_em & {\scriptsize texto(40)} \\ {\scriptsize }ultima\_atividade\_em & {\scriptsize texto(40)} \\ {\scriptsize }encerrada\_em & {\scriptsize texto(40)?} \\ {\scriptsize }motivo\_encerramento & {\scriptsize texto(40)}\end{tabular}};
\node[tabela] (anotacoes) at (6.5,0) {\textbf{anotacoes}\nodepart{second}\begin{tabular}{@{}l@{\hspace{4pt}}l@{}}{\scriptsize PK }\underline{id} & {\scriptsize texto(32)} \\ {\scriptsize FK }caso\_id & {\scriptsize texto(32)} \\ {\scriptsize }alvo\_tipo & {\scriptsize texto(10)} \\ {\scriptsize }alvo\_id & {\scriptsize texto(32)} \\ {\scriptsize }autor & {\scriptsize texto(100)} \\ {\scriptsize }criado\_em & {\scriptsize texto(40)} \\ {\scriptsize }texto & {\scriptsize texto}\end{tabular}};
\node[tabela] (tarefas) at (6.5,-6.6) {\textbf{tarefas}\nodepart{second}\begin{tabular}{@{}l@{\hspace{4pt}}l@{}}{\scriptsize PK }\underline{id} & {\scriptsize texto(32)} \\ {\scriptsize FK }caso\_id & {\scriptsize texto(32)} \\ {\scriptsize }titulo & {\scriptsize texto(300)} \\ {\scriptsize }descricao & {\scriptsize texto} \\ {\scriptsize }responsavel & {\scriptsize texto(100)?} \\ {\scriptsize }situacao & {\scriptsize texto(15)} \\ {\scriptsize }prazo & {\scriptsize texto(10)?} \\ {\scriptsize }criado\_por & {\scriptsize texto(100)} \\ {\scriptsize }criado\_em & {\scriptsize texto(40)} \\ {\scriptsize }atualizado\_em & {\scriptsize texto(40)}\end{tabular}};
\node[tabela] (tarefaeventos) at (6.5,-13.6) {\textbf{tarefa\_eventos}\nodepart{second}\begin{tabular}{@{}l@{\hspace{4pt}}l@{}}{\scriptsize PK }\underline{id} & {\scriptsize inteiro} \\ {\scriptsize FK }tarefa\_id & {\scriptsize texto(32)} \\ {\scriptsize FK }caso\_id & {\scriptsize texto(32)} \\ {\scriptsize }usuario & {\scriptsize texto(100)} \\ {\scriptsize }quando & {\scriptsize texto(40)} \\ {\scriptsize }alteracoes & {\scriptsize JSON}\end{tabular}};
\node[tabela] (auditoria) at (0,-14.0) {\textbf{auditoria}\nodepart{second}\begin{tabular}{@{}l@{\hspace{4pt}}l@{}}{\scriptsize PK }\underline{seq} & {\scriptsize inteiro} \\ {\scriptsize }versao & {\scriptsize inteiro} \\ {\scriptsize }quando & {\scriptsize texto(40)} \\ {\scriptsize }usuario & {\scriptsize texto(100)} \\ {\scriptsize }acao & {\scriptsize texto(60)} \\ {\scriptsize }caso\_id & {\scriptsize texto(32)?} \\ {\scriptsize }justificativa & {\scriptsize texto?} \\ {\scriptsize }detalhes & {\scriptsize JSON} \\ {\scriptsize }hash\_anterior & {\scriptsize texto(64)} \\ {\scriptsize }hash & {\scriptsize texto(64)}\end{tabular}};
\draw[fk] (anotacoes) to[] node[mult,pos=0.2]{0..*} node[mult,pos=0.8]{1} (casos);
\draw[fk] (casomembros) to[] node[mult,pos=0.2]{0..*} node[mult,pos=0.8]{1} (casos);
\draw[fk] (sessoes) to[] node[mult,pos=0.2]{0..*} node[mult,pos=0.8]{1} (usuarios);
\draw[fk] (tarefas) to[] node[mult,pos=0.2]{0..*} node[mult,pos=0.8]{1} (casos);
\draw[fk] (tarefaeventos) to[] node[mult,pos=0.2]{0..*} node[mult,pos=0.8]{1} (tarefas);
\draw[fk] (tarefaeventos) to[] node[mult,pos=0.2]{0..*} node[mult,pos=0.8]{1} (casos);
\draw[logica] (casomembros) -- node[mult,pos=0.5]{usuario} (usuarios);
\draw[logica] (tarefas) -- node[mult,pos=0.5]{responsavel} (usuarios);
\draw[logica] (auditoria) -- node[mult,pos=0.5]{caso\_id} (casos);
\end{tikzpicture}}
\fonte{Gerada pelo autor a partir do modelo SQLAlchemy do back-end (\texttt{buscador/modelos.py}). PK: chave primária; FK: chave estrangeira; ``?'': aceita nulo. Linha cheia: chave estrangeira; linha tracejada: referência sem chave estrangeira. Multiplicidade junto de cada ponta. Tabelas mostradas só com as chaves: colunas completas no Apêndice~\ref{ap:dicionario}.}
\end{figure}


% Gerado por docs/sistema/gera_modelo_tex.py a partir de docs/sistema/modelo.json. Não editar à mão.

\begin{small}
\setlength{\tabcolsep}{3pt}
\begin{longtable}{@{}p{2.9cm}p{3.5cm}p{2.9cm}p{1.5cm}p{1.0cm}p{3.0cm}@{}}
\caption{Relações do modelo de dados e suas cardinalidades}\rotulounico{tab:relacoes}\\
\hline
\textbf{Tabela ``1''} & \textbf{Tabela relacionada} & \textbf{Coluna} & \textbf{Multipl.} & \textbf{Card.} & \textbf{Tipo} \\ \hline\endhead
\nometeste{alertas} & \nometeste{alerta_decisoes} & \nometeste{alerta_id} & 1 : 0..* & 1:N & chave estrangeira \\
\nometeste{capturas} & \nometeste{arestas} & \nometeste{captura_id} & 1 : 0..* & 1:N & chave estrangeira \\
\nometeste{capturas} & \nometeste{entidade_fontes} & \nometeste{captura_id} & 1 : 0..* & 1:N & chave estrangeira \\
\nometeste{capturas} & \nometeste{entidades} & \nometeste{captura_id} & 1 : 0..* & 1:N & chave estrangeira \\
\nometeste{capturas} & \nometeste{extracoes_ia} & \nometeste{captura_id} & 1 : 0..* & 1:N & chave estrangeira \\
\nometeste{capturas} & \nometeste{sugestoes_ia} & \nometeste{captura_id} & 1 : 0..* & 1:N & chave estrangeira \\
\nometeste{casos} & \nometeste{alerta_decisoes} & \nometeste{caso_id} & 1 : 0..* & 1:N & chave estrangeira \\
\nometeste{casos} & \nometeste{alertas} & \nometeste{caso_id} & 1 : 0..* & 1:N & chave estrangeira \\
\nometeste{casos} & \nometeste{anotacoes} & \nometeste{caso_id} & 1 : 0..* & 1:N & chave estrangeira \\
\nometeste{casos} & \nometeste{arestas} & \nometeste{caso_id} & 1 : 0..* & 1:N & chave estrangeira \\
\nometeste{casos} & \nometeste{avisos} & \nometeste{caso_id} & 1 : 0..* & 1:N & chave estrangeira \\
\nometeste{casos} & \nometeste{capturas} & \nometeste{caso_id} & 1 : 0..* & 1:N & chave estrangeira \\
\nometeste{casos} & \nometeste{caso_membros} & \nometeste{caso_id} & 1 : 0..* & 1:N & chave estrangeira \\
\nometeste{casos} & \nometeste{coletas} & \nometeste{caso_id} & 1 : 0..* & 1:N & chave estrangeira \\
\nometeste{casos} & \nometeste{entidades} & \nometeste{caso_id} & 1 : 0..* & 1:N & chave estrangeira \\
\nometeste{casos} & \nometeste{extracoes_ia} & \nometeste{caso_id} & 1 : 0..* & 1:N & chave estrangeira \\
\nometeste{casos} & \nometeste{monitoramento_execucoes} & \nometeste{caso_id} & 1 : 0..* & 1:N & chave estrangeira \\
\nometeste{casos} & \nometeste{monitoramentos} & \nometeste{caso_id} & 1 : 0..* & 1:N & chave estrangeira \\
\nometeste{casos} & \nometeste{sugestoes} & \nometeste{caso_id} & 1 : 0..* & 1:N & chave estrangeira \\
\nometeste{casos} & \nometeste{sugestoes_ia} & \nometeste{caso_id} & 1 : 0..* & 1:N & chave estrangeira \\
\nometeste{casos} & \nometeste{tarefa_eventos} & \nometeste{caso_id} & 1 : 0..* & 1:N & chave estrangeira \\
\nometeste{casos} & \nometeste{tarefas} & \nometeste{caso_id} & 1 : 0..* & 1:N & chave estrangeira \\
\nometeste{coletas} & \nometeste{capturas} & \nometeste{coleta_id} & 0..1 : 0..* & 1:N & chave estrangeira \\
\nometeste{coletas} & \nometeste{coletas} & \nometeste{recoleta_de} & 0..1 : 0..* & 1:N & chave estrangeira \\
\nometeste{entidades} & \nometeste{arestas} & \nometeste{destino_id} & 1 : 0..* & 1:N & chave estrangeira \\
\nometeste{entidades} & \nometeste{arestas} & \nometeste{origem_id} & 1 : 0..* & 1:N & chave estrangeira \\
\nometeste{entidades} & \nometeste{avisos} & \nometeste{entidade_id} & 1 : 0..* & 1:N & chave estrangeira \\
\nometeste{entidades} & \nometeste{entidade_fontes} & \nometeste{entidade_id} & 1 : 0..* & 1:N & chave estrangeira \\
\nometeste{entidades} & \nometeste{monitoramentos} & \nometeste{entidade_id} & 1 : 0..* & 1:N & chave estrangeira \\
\nometeste{entidades} & \nometeste{sugestoes} & \nometeste{candidata_id} & 1 : 0..* & 1:N & chave estrangeira \\
\nometeste{entidades} & \nometeste{sugestoes} & \nometeste{entidade_id} & 1 : 0..* & 1:N & chave estrangeira \\
\nometeste{entidades} & \nometeste{sugestoes_ia} & \nometeste{entidade_id} & 0..1 : 0..* & 1:N & chave estrangeira \\
\nometeste{extracoes_ia} & \nometeste{sugestoes_ia} & \nometeste{extracao_id} & 1 : 0..* & 1:N & chave estrangeira \\
\nometeste{ij_decisoes} & \nometeste{ij_importacoes} & \nometeste{decisao_id} & 1 : 0..* & 1:N & chave estrangeira \\
\nometeste{ij_decisoes} & \nometeste{ij_registros} & \nometeste{decisao_id} & 1 : 0..* & 1:N & chave estrangeira \\
\nometeste{ij_importacoes} & \nometeste{ij_arquivos} & \nometeste{importacao_id} & 1 : 0..* & 1:N & chave estrangeira \\
\nometeste{ij_importacoes} & \nometeste{ij_registros} & \nometeste{importacao_id} & 1 : 0..* & 1:N & chave estrangeira \\
\nometeste{monitoramento_execucoes} & \nometeste{avisos} & \nometeste{execucao_id} & 1 : 0..* & 1:N & chave estrangeira \\
\nometeste{monitoramentos} & \nometeste{avisos} & \nometeste{monitoramento_id} & 1 : 0..* & 1:N & chave estrangeira \\
\nometeste{monitoramentos} & \nometeste{monitoramento_execucoes} & \nometeste{monitoramento_id} & 1 : 0..* & 1:N & chave estrangeira \\
\nometeste{tarefas} & \nometeste{tarefa_eventos} & \nometeste{tarefa_id} & 1 : 0..* & 1:N & chave estrangeira \\
\nometeste{usuarios} & \nometeste{sessoes} & \nometeste{conta_id} & 1 : 0..* & 1:N & chave estrangeira \\
\nometeste{entidades} & \nometeste{alertas} & \nometeste{entidades} & -- & N:N & lista JSON de IDs de entidades \\
\nometeste{capturas} & \nometeste{alertas} & \nometeste{evidencias} & -- & N:N & lista JSON de evidências com o ID da captura \\
\nometeste{capturas} & \nometeste{avisos} & \nometeste{captura_ids} & -- & N:N & lista JSON de IDs de capturas \\
\nometeste{coletas} & \nometeste{monitoramento_execucoes} & \nometeste{coletas} & -- & N:N & lista JSON com as recoletas feitas \\
\nometeste{entidades} & \nometeste{anotacoes} & \nometeste{alvo_id} & -- & 1:N & alvo polimórfico (alvo\_tipo = entidade) \\
\nometeste{alertas} & \nometeste{anotacoes} & \nometeste{alvo_id} & -- & 1:N & alvo polimórfico (alvo\_tipo = alerta) \\
\nometeste{usuarios} & \nometeste{caso_membros} & \nometeste{usuario} & -- & 1:N & login; o usuário de token de desenvolvimento não tem linha \\
\nometeste{usuarios} & \nometeste{tarefas} & \nometeste{responsavel} & -- & 1:N & login do responsável, opcional \\
\nometeste{casos} & \nometeste{auditoria} & \nometeste{caso_id} & -- & 1:N & sem chave estrangeira: a trilha só de acréscimo não depende do caso \\
\nometeste{casos} & \nometeste{ij_decisoes} & \nometeste{caso_id} & -- & 1:N & sem chave estrangeira: base separada (RF10) \\
\nometeste{casos} & \nometeste{ij_importacoes} & \nometeste{caso_id} & -- & 1:N & sem chave estrangeira: base separada (RF10) \\
\nometeste{casos} & \nometeste{ij_registros} & \nometeste{caso_id} & -- & 1:N & sem chave estrangeira: base separada (RF10) \\
\hline
\multicolumn{6}{p{15cm}}{\footnotesize Fonte: gerada pelo autor a partir de \texttt{buscador/modelos.py}. Multiplicidade ``1 : 0..*'' significa que cada linha da tabela relacionada aponta para exatamente uma linha da tabela ``1'', que pode ter nenhuma ou várias. Referências sem chave estrangeira foram conferidas no código.}\\
\end{longtable}
\end{small}


Algumas referências não usam chave estrangeira, por decisão de projeto:
\begin{itemize}
	\item \textbf{Listas em JSON:} as entidades e as evidências de um alerta,
	as capturas de um aviso e as recoletas de uma execução do monitoramento.
	Cada lista é gravada junto com o registro que a produziu e não muda depois.
	\item \textbf{Alvo polimórfico:} uma anotação aponta para uma entidade ou
	para um alerta, conforme \texttt{alvo\_tipo}. A API confere se o alvo é do
	mesmo caso.
	\item \textbf{Login:} membros e responsáveis de tarefas guardam o login,
	porque o usuário de token de desenvolvimento não tem linha em
	\texttt{usuarios}.
	\item \textbf{Auditoria:} a trilha guarda o caso sem chave estrangeira,
	para não depender de nenhuma outra tabela.
\end{itemize}

% ===========================================================================
\section{Fluxos principais}
\label{sec:desenv-fluxos}

Três fluxos sustentam o diferencial do protótipo, a produção de prova para o
inquérito (Seção~\ref{sec:requisitos-mercado}). A
Figura~\ref{fig:uml-seq-coleta} mostra a coleta com cadeia de custódia; a
Figura~\ref{fig:uml-seq-alerta}, a execução das tipologias e a decisão sobre um
alerta; e a Figura~\ref{fig:uml-seq-verificacao}, a verificação de integridade de
uma captura e da trilha de auditoria. As Figuras~\ref{fig:uml-estado-alerta}
e~\ref{fig:uml-estado-caso} mostram os estados do alerta e do caso.

% Diagrama de sequência (UML 2): coleta com custódia.
% Elementos e fontes no código: docs/sistema/diagramas-fontes.md.
% Coordenadas geradas por script; para mudar, edite as coordenadas à mão.
\begin{figure}[p]
\centering
\caption{Diagrama de sequência de uma coleta com registro na cadeia de custódia}
\label{fig:uml-seq-coleta}
\resizebox{0.95\textwidth}{!}{%
\begin{tikzpicture}[
  font=\footnotesize,
  vida/.style={draw, fill=white, align=center, inner sep=3pt, font=\footnotesize, minimum height=0.9cm},
  rotvida/.style={align=center, font=\footnotesize, inner sep=1pt},
  linhavida/.style={dashed, gray!70!black},
  msgsinc/.style={-{Latex[length=5pt]}},
  msgresp/.style={-{Straight Barb[length=5pt]}, dashed},
  rotmsg/.style={font=\footnotesize, align=center, inner sep=1pt, fill=white},
  frag/.style={draw, thin},
  rotfrag/.style={font=\footnotesize\bfseries, inner sep=2.5pt},
  rotguarda/.style={font=\footnotesize, inner sep=1pt, fill=white},
  notaseq/.style={align=left, font=\scriptsize, inner sep=3pt},
  ator/.pic={
    \draw (0,0.32) circle (0.12);
    \draw (0,0.2) -- (0,-0.12) (-0.22,0.1) -- (0.22,0.1)
          (0,-0.12) -- (-0.17,-0.42) (0,-0.12) -- (0.17,-0.42);}
]
\providecommand{\dobra}[1]{}
\renewcommand{\dobra}[1]{%
  \draw (#1.north west) -- ([xshift=-5pt]#1.north east) -- ([yshift=-5pt]#1.north east)
        -- (#1.south east) -- (#1.south west) -- cycle;
  \draw ([xshift=-5pt]#1.north east) |- ([yshift=-5pt]#1.north east);}
\pic at (0.00,0.55) {ator};
\node[rotvida, anchor=north] at (0.00,-0.05) {Investigador};
\draw[linhavida] (0.00,-0.75) -- (0.00,-29.94);
\node[vida] at (2.50,-0.15) {Front-end\\{\scriptsize\texttt{AbaColeta}}};
\draw[linhavida] (2.50,-0.75) -- (2.50,-29.94);
\node[vida] at (5.40,-0.15) {API\\{\scriptsize\texttt{rotas\_casos}, \texttt{coleta}}};
\draw[linhavida] (5.40,-0.75) -- (5.40,-29.94);
\node[vida] at (9.40,-0.15) {Auditoria\\{\scriptsize\texttt{auditoria}}};
\draw[linhavida] (9.40,-0.75) -- (9.40,-29.94);
\node[vida] at (11.70,-0.15) {Resolução\\{\scriptsize\texttt{Resolvedor}}};
\draw[linhavida] (11.70,-0.75) -- (11.70,-29.94);
\node[vida] at (14.00,-0.15) {Conector};
\draw[linhavida] (14.00,-0.75) -- (14.00,-29.94);
\node[vida] at (16.40,-0.15) {Fonte\\{\scriptsize PNCP, Portal}\\{\scriptsize ou arquivo}};
\draw[linhavida] (16.40,-0.75) -- (16.40,-29.94);
\node[vida] at (18.60,-0.15) {Custódia\\{\scriptsize\texttt{Custodia}}};
\draw[linhavida] (18.60,-0.75) -- (18.60,-29.94);
\draw[msgsinc] (0.00,-1.94) -- (2.50,-1.94);
\node[rotmsg, above=1pt] at (1.25,-1.94) {fonte e parâmetros\\(justificativa já na barra)};
\draw[msgsinc] (2.50,-2.98) -- (5.40,-2.98);
\node[rotmsg, above=1pt] at (3.95,-2.98) {\texttt{POST /casos/\{id\}/coletas}, via nginx \texttt{/api/}\\\texttt{Authorization: Bearer}, \texttt{X-Justificativa}};
\draw[msgsinc] (5.40,-3.36) -- ++(0.35,0) -- ++(0,-1.14) -- (5.40,-4.50);
\node[rotmsg, anchor=north west, align=left] at (5.85,-3.28) {sessão ativa e justificativa\\(\texttt{usuario\_atual},\\ \texttt{justificativa\_obrigatoria})};
\draw[msgsinc] (5.40,-4.88) -- ++(0.35,0) -- ++(0,-1.14) -- (5.40,-6.02);
\node[rotmsg, anchor=north west, align=left] at (5.85,-4.80) {\texttt{caso\_aberto}: papel $\geq$\\investigador e prazo\\de retenção em vigor};
\draw[msgsinc] (5.40,-7.61) -- (9.40,-7.61);
\node[rotmsg, above=1pt] at (7.40,-7.61) {registrar \texttt{caso.acesso\_negado}\\ou \texttt{caso.retencao\_recusa}};
\draw[msgresp] (5.40,-8.27) -- (2.50,-8.27);
\node[rotmsg, above=1pt] at (3.95,-8.27) {403};
\draw[frag] (-0.45,-6.27) rectangle (22.35,-8.52);
\node[rotfrag, anchor=north west] (fr) at (-0.45,-6.27) {break};
\draw[frag] (fr.south west) -- ([xshift=-0.12cm]fr.south east) -- ([yshift=0.12cm]fr.south east) -- (fr.north east);
\node[rotguarda, anchor=north west] at (0.75,-6.30) {[não é membro, papel insuficiente ou prazo vencido]};
\draw[msgsinc] (5.40,-9.66) -- (9.40,-9.66);
\node[rotmsg, above=1pt] at (7.40,-9.66) {linha em \texttt{coletas};\\registrar \texttt{coleta.iniciar}; commit; âncora};
\draw[msgsinc] (5.40,-10.32) -- (14.00,-10.32);
\node[rotmsg, above=1pt] at (9.70,-10.32) {\texttt{criar\_conector}; \texttt{coletar(parâmetros, registrar)}};
\draw[msgsinc] (14.00,-13.01) -- (16.40,-13.01);
\node[rotmsg, above=1pt] at (15.20,-13.01) {consulta HTTPS\\ou leitura};
\draw[msgresp] (16.40,-13.67) -- (14.00,-13.67);
\node[rotmsg, above=1pt] at (15.20,-13.67) {carga bruta};
\draw[frag] (9.10,-11.67) rectangle (21.90,-13.92);
\node[rotfrag, anchor=north west] (fr) at (9.10,-11.67) {loop};
\draw[frag] (fr.south west) -- ([xshift=-0.12cm]fr.south east) -- ([yshift=0.12cm]fr.south east) -- (fr.north east);
\node[rotguarda, anchor=north west] at (10.15,-11.70) {[Retentativa: até $N$ tentativas, espera $e \times n$ após falha transitória]};
\draw[msgsinc] (14.00,-14.68) -- (18.60,-14.68);
\node[rotmsg, above=1pt] at (16.30,-14.68) {\texttt{registrar(carga)}};
\draw[msgsinc] (18.60,-15.06) -- ++(0.35,0) -- ++(0,-0.76) -- (18.60,-15.82);
\node[rotmsg, anchor=north west, align=left] at (19.05,-14.98) {SHA-256 do\\conteúdo original};
\draw[msgsinc] (18.60,-16.20) -- ++(0.35,0) -- ++(0,-1.14) -- (18.60,-17.34);
\node[rotmsg, anchor=north west, align=left] at (19.05,-16.12) {AES-256-GCM\\(subchave HKDF,\\AAD = SHA-256)};
\draw[msgsinc] (18.60,-17.72) -- ++(0.35,0) -- ++(0,-1.90) -- (18.60,-19.62);
\node[rotmsg, anchor=north west, align=left] at (19.05,-17.64) {grava \texttt{.gcm}\\(temporário e\\\texttt{os.replace});\\\texttt{INSERT} em\\\texttt{capturas}};
\draw[msgresp] (18.60,-20.28) -- (14.00,-20.28);
\node[rotmsg, above=1pt] at (16.30,-20.28) {captura};
\draw[frag] (8.95,-11.12) rectangle (22.05,-20.53);
\node[rotfrag, anchor=north west] (fr) at (8.95,-11.12) {loop};
\draw[frag] (fr.south west) -- ([xshift=-0.12cm]fr.south east) -- ([yshift=0.12cm]fr.south east) -- (fr.north east);
\node[rotguarda, anchor=north west] at (10.00,-11.15) {[cada página ou arquivo]};
\draw[msgresp] (14.00,-21.29) -- (5.40,-21.29);
\node[rotmsg, above=1pt] at (9.70,-21.29) {registros e capturas};
\draw[msgsinc] (5.40,-21.95) -- (11.70,-21.95);
\node[rotmsg, above=1pt] at (8.55,-21.95) {\texttt{incorporar(registros)} em SAVEPOINT};
\draw[msgresp] (11.70,-22.61) -- (5.40,-22.61);
\node[rotmsg, above=1pt] at (8.55,-22.61) {entidades, arestas e sugestões};
\draw[msgsinc] (5.40,-23.65) -- (9.40,-23.65);
\node[rotmsg, above=1pt] at (7.40,-23.65) {registrar \texttt{coleta.executar}\\(HMAC encadeado); commit; âncora};
\draw[msgresp] (5.40,-24.31) -- (2.50,-24.31);
\node[rotmsg, above=1pt] at (3.95,-24.31) {200: capturas, entidades novas, tentativas};
\draw[msgresp] (2.50,-24.97) -- (0.00,-24.97);
\node[rotmsg, above=1pt] at (1.25,-24.97) {resultado};
\draw[msgsinc] (5.40,-26.18) -- (9.40,-26.18);
\node[rotmsg, above=1pt] at (7.40,-26.18) {registrar \texttt{coleta.nao\_configurada}};
\draw[msgresp] (5.40,-26.84) -- (2.50,-26.84);
\node[rotmsg, above=1pt] at (3.95,-26.84) {200, situação \texttt{nao\_configurado}};
\draw[msgsinc] (5.40,-28.43) -- (9.40,-28.43);
\node[rotmsg, above=1pt] at (7.40,-28.43) {registrar \texttt{coleta.falhou};\\capturas já guardadas confirmadas};
\draw[msgresp] (5.40,-29.09) -- (2.50,-29.09);
\node[rotmsg, above=1pt] at (3.95,-29.09) {422 ou 500};
\draw[frag] (-0.30,-10.57) rectangle (22.20,-29.34);
\node[rotfrag, anchor=north west] (fr) at (-0.30,-10.57) {alt};
\draw[frag] (fr.south west) -- ([xshift=-0.12cm]fr.south east) -- ([yshift=0.12cm]fr.south east) -- (fr.north east);
\draw[frag, dashed] (-0.30,-25.22) -- (22.20,-25.22);
\draw[frag, dashed] (-0.30,-27.09) -- (22.20,-27.09);
\node[rotguarda, anchor=north west] at (0.60,-10.60) {[sem erro]};
\node[rotguarda, anchor=north west] at (-0.18,-25.25) {[fonte não configurada (\texttt{FonteNaoConfigurada})]};
\node[rotguarda, anchor=north west] at (-0.18,-27.12) {[outra falha (conector, resolução ou gravação)]};
\end{tikzpicture}}
\fonte{Elaborada pelo autor a partir de \texttt{buscador/api/rotas\_casos.py} (\texttt{coletar}, \texttt{\_coletar\_registrando}), \texttt{buscador/api/deps.py}, \texttt{buscador/coleta.py}, \texttt{buscador/conectores/base.py} (\texttt{Retentativa}), \texttt{buscador/custodia/servico.py}, \texttt{buscador/cripto.py}, \texttt{buscador/auditoria/servico.py} e, no front-end, \texttt{src/telas/abas/AbaColeta.tsx} e \texttt{src/api/clienteHttp.ts}. ``Commit; âncora'': confirmação da transação seguida da cópia do topo da trilha para a âncora externa.}
\end{figure}

% Diagrama de sequência (UML 2): tipologias, alertas e decisão.
% Elementos e fontes no código: docs/sistema/diagramas-fontes.md.
% Coordenadas geradas por script; para mudar, edite as coordenadas à mão.
\begin{figure}[htbp]
\centering
\caption{Diagrama de sequência da execução das tipologias e da decisão sobre um alerta}
\label{fig:uml-seq-alerta}
\resizebox{0.95\textwidth}{!}{%
\begin{tikzpicture}[
  font=\footnotesize,
  vida/.style={draw, fill=white, align=center, inner sep=3pt, font=\footnotesize, minimum height=0.9cm},
  rotvida/.style={align=center, font=\footnotesize, inner sep=1pt},
  linhavida/.style={dashed, gray!70!black},
  msgsinc/.style={-{Latex[length=5pt]}},
  msgresp/.style={-{Straight Barb[length=5pt]}, dashed},
  rotmsg/.style={font=\footnotesize, align=center, inner sep=1pt, fill=white},
  frag/.style={draw, thin},
  rotfrag/.style={font=\footnotesize\bfseries, inner sep=2.5pt},
  rotguarda/.style={font=\footnotesize, inner sep=1pt, fill=white},
  notaseq/.style={align=left, font=\scriptsize, inner sep=3pt},
  ator/.pic={
    \draw (0,0.32) circle (0.12);
    \draw (0,0.2) -- (0,-0.12) (-0.22,0.1) -- (0.22,0.1)
          (0,-0.12) -- (-0.17,-0.42) (0,-0.12) -- (0.17,-0.42);}
]
\providecommand{\dobra}[1]{}
\renewcommand{\dobra}[1]{%
  \draw (#1.north west) -- ([xshift=-5pt]#1.north east) -- ([yshift=-5pt]#1.north east)
        -- (#1.south east) -- (#1.south west) -- cycle;
  \draw ([xshift=-5pt]#1.north east) |- ([yshift=-5pt]#1.north east);}
\pic at (0.00,0.55) {ator};
\node[rotvida, anchor=north] at (0.00,-0.05) {Investigador};
\draw[linhavida] (0.00,-0.75) -- (0.00,-23.88);
\node[vida] at (2.60,-0.15) {Front-end\\{\scriptsize\texttt{AbaAlertas}}};
\draw[linhavida] (2.60,-0.75) -- (2.60,-23.88);
\node[vida] at (5.60,-0.15) {API\\{\scriptsize\texttt{rotas\_casos}}};
\draw[linhavida] (5.60,-0.75) -- (5.60,-23.88);
\node[vida] at (9.60,-0.15) {Tipologias\\{\scriptsize\texttt{tipologias.motor}}};
\draw[linhavida] (9.60,-0.75) -- (9.60,-23.88);
\node[vida] at (13.60,-0.15) {Banco\\{\scriptsize PostgreSQL}};
\draw[linhavida] (13.60,-0.75) -- (13.60,-23.88);
\node[vida] at (16.40,-0.15) {Auditoria\\{\scriptsize\texttt{auditoria}}};
\draw[linhavida] (16.40,-0.75) -- (16.40,-23.88);
\draw[msgsinc] (0.00,-1.94) -- (2.60,-1.94);
\node[rotmsg, above=1pt] at (1.30,-1.94) {executar\\tipologias};
\draw[msgsinc] (2.60,-2.98) -- (5.60,-2.98);
\node[rotmsg, above=1pt] at (4.10,-2.98) {\texttt{POST /casos/\{id\}/alertas/executar}\\\texttt{X-Justificativa}};
\draw[msgsinc] (5.60,-3.36) -- ++(0.35,0) -- ++(0,-1.52) -- (5.60,-4.88);
\node[rotmsg, anchor=north west, align=left] at (6.05,-3.28) {\texttt{caso\_aberto}: papel $\geq$\\investigador e prazo em vigor;\\parâmetros efetivos\\(configuração e sobrescrita)};
\draw[msgsinc] (5.60,-5.92) -- (9.60,-5.92);
\node[rotmsg, above=1pt] at (7.60,-5.92) {\texttt{executar(caso, parâmetros,}\\\texttt{ler\_captura)}};
\draw[msgsinc] (9.60,-6.96) -- (13.60,-6.96);
\node[rotmsg, above=1pt] at (11.60,-6.96) {trava o caso (\texttt{FOR UPDATE});\\lê entidades, arestas e capturas};
\draw[msgresp] (13.60,-7.62) -- (9.60,-7.62);
\node[rotmsg, above=1pt] at (11.60,-7.62) {contexto do caso};
\draw[msgsinc] (9.60,-8.55) -- ++(0.35,0) -- ++(0,-1.52) -- (9.60,-10.07);
\node[rotmsg, anchor=north west, align=left] at (10.05,-8.47) {regra $\to$ alerta detectado:\\motivo com os números reais,\\evidências (captura, fonte,\\SHA-256) e chave estável};
\draw[frag] (5.30,-7.87) rectangle (14.80,-10.32);
\node[rotfrag, anchor=north west] (fr) at (5.30,-7.87) {loop};
\draw[frag] (fr.south west) -- ([xshift=-0.12cm]fr.south east) -- ([yshift=0.12cm]fr.south east) -- (fr.north east);
\node[rotguarda, anchor=north west] at (6.35,-7.90) {[cada tipologia pedida (T1 a T10)]};
\draw[msgsinc] (9.60,-11.84) -- (13.60,-11.84);
\node[rotmsg, above=1pt] at (11.60,-11.84) {reconcilia pela chave: atualiza ou cria;\\apaga pendente sem decisão;\\decidido não detectado: \texttt{ativo = falso}};
\draw[msgresp] (9.60,-12.50) -- (5.60,-12.50);
\node[rotmsg, above=1pt] at (7.60,-12.50) {alertas por severidade};
\draw[msgsinc] (5.60,-13.54) -- (16.40,-13.54);
\node[rotmsg, above=1pt] at (11.00,-13.54) {registrar \texttt{alerta.executar};\\commit; âncora};
\draw[msgresp] (5.60,-14.20) -- (2.60,-14.20);
\node[rotmsg, above=1pt] at (4.10,-14.20) {alertas com motivo e evidências};
\draw[msgresp] (2.60,-14.86) -- (0.00,-14.86);
\node[rotmsg, above=1pt] at (1.30,-14.86) {lista};
\draw[msgsinc] (0.00,-16.10) -- (2.60,-16.10);
\node[rotmsg, above=1pt] at (1.30,-16.10) {confirma ou descarta,\\com justificativa};
\draw[msgsinc] (2.60,-17.52) -- (5.60,-17.52);
\node[rotmsg, above=1pt] at (4.10,-17.52) {\texttt{POST .../alertas/\{aid\}/decisao}\\corpo: \texttt{decisao}, \texttt{justificativa};\\cabeçalho: \texttt{X-Justificativa}};
\draw[msgsinc] (5.60,-17.90) -- ++(0.35,0) -- ++(0,-0.76) -- (5.60,-18.66);
\node[rotmsg, anchor=north west, align=left] at (6.05,-17.82) {\texttt{caso\_aberto}: papel $\geq$\\investigador e prazo em vigor};
\draw[msgsinc] (5.60,-19.32) -- (13.60,-19.32);
\node[rotmsg, above=1pt] at (9.60,-19.32) {\texttt{SELECT} do alerta \texttt{FOR UPDATE}};
\draw[msgresp] (13.60,-19.98) -- (5.60,-19.98);
\node[rotmsg, above=1pt] at (9.60,-19.98) {alerta e situação anterior};
\draw[msgsinc] (5.60,-21.02) -- (13.60,-21.02);
\node[rotmsg, above=1pt] at (9.60,-21.02) {\texttt{INSERT} em \texttt{alerta\_decisoes};\\\texttt{UPDATE} de \texttt{alertas} (situação e última decisão)};
\draw[msgsinc] (5.60,-22.06) -- (16.40,-22.06);
\node[rotmsg, above=1pt] at (11.00,-22.06) {registrar \texttt{alerta.decidir}: decisão, justificativa,\\situação anterior e \texttt{revisao}; commit; âncora};
\draw[msgresp] (5.60,-22.72) -- (2.60,-22.72);
\node[rotmsg, above=1pt] at (4.10,-22.72) {alerta com o histórico de decisões};
\draw[msgresp] (2.60,-23.38) -- (0.00,-23.38);
\node[rotmsg, above=1pt] at (1.30,-23.38) {situação atualizada};
\end{tikzpicture}}
\fonte{Elaborada pelo autor a partir de \texttt{buscador/api/rotas\_casos.py} (\texttt{executar\_alertas}, \texttt{decidir\_alerta}, \texttt{\_obter\_alerta}), \texttt{buscador/tipologias/motor.py} (\texttt{executar}), \texttt{buscador/modelos.py} (\texttt{Alerta}, \texttt{AlertaDecisao}), \texttt{buscador/api/deps.py} e, no front-end, \texttt{src/telas/abas/AbaAlertas.tsx}. A recusa por papel ou por prazo vencido segue o fluxo da Figura~\ref{fig:uml-seq-coleta}.}
\end{figure}

% Diagrama de sequência (UML 2): verificação de integridade.
% Elementos e fontes no código: docs/sistema/diagramas-fontes.md.
% Coordenadas geradas por script; para mudar, edite as coordenadas à mão.
\begin{figure}[p]
\centering
\caption{Diagrama de sequência da verificação de integridade de uma captura e da cadeia da auditoria}
\label{fig:uml-seq-verificacao}
\resizebox{0.95\textwidth}{!}{%
\begin{tikzpicture}[
  font=\footnotesize,
  vida/.style={draw, fill=white, align=center, inner sep=3pt, font=\footnotesize, minimum height=0.9cm},
  rotvida/.style={align=center, font=\footnotesize, inner sep=1pt},
  linhavida/.style={dashed, gray!70!black},
  msgsinc/.style={-{Latex[length=5pt]}},
  msgresp/.style={-{Straight Barb[length=5pt]}, dashed},
  rotmsg/.style={font=\footnotesize, align=center, inner sep=1pt, fill=white},
  frag/.style={draw, thin},
  rotfrag/.style={font=\footnotesize\bfseries, inner sep=2.5pt},
  rotguarda/.style={font=\footnotesize, inner sep=1pt, fill=white},
  notaseq/.style={align=left, font=\scriptsize, inner sep=3pt},
  ator/.pic={
    \draw (0,0.32) circle (0.12);
    \draw (0,0.2) -- (0,-0.12) (-0.22,0.1) -- (0.22,0.1)
          (0,-0.12) -- (-0.17,-0.42) (0,-0.12) -- (0.17,-0.42);}
]
\providecommand{\dobra}[1]{}
\renewcommand{\dobra}[1]{%
  \draw (#1.north west) -- ([xshift=-5pt]#1.north east) -- ([yshift=-5pt]#1.north east)
        -- (#1.south east) -- (#1.south west) -- cycle;
  \draw ([xshift=-5pt]#1.north east) |- ([yshift=-5pt]#1.north east);}
\pic at (0.00,0.55) {ator};
\node[rotvida, anchor=north] at (0.00,-0.05) {Usuário\\autenticado};
\draw[linhavida] (0.00,-0.75) -- (0.00,-29.30);
\node[vida] at (2.60,-0.15) {Front-end\\{\scriptsize\texttt{CartaoAlerta},}\\{\scriptsize\texttt{AbaAuditoria}}};
\draw[linhavida] (2.60,-0.75) -- (2.60,-29.30);
\node[vida] at (5.80,-0.15) {API\\{\scriptsize\texttt{rotas\_capturas},}\\{\scriptsize\texttt{rotas\_auditoria}}};
\draw[linhavida] (5.80,-0.75) -- (5.80,-29.30);
\node[vida] at (10.00,-0.15) {Custódia\\{\scriptsize\texttt{Custodia}}};
\draw[linhavida] (10.00,-0.75) -- (10.00,-29.30);
\node[vida] at (12.80,-0.15) {Banco\\{\scriptsize\texttt{auditoria}}};
\draw[linhavida] (12.80,-0.75) -- (12.80,-29.30);
\node[vida] at (15.00,-0.15) {Âncora\\{\scriptsize\texttt{auditoria.json}}};
\draw[linhavida] (15.00,-0.75) -- (15.00,-29.30);
\node[vida] at (17.20,-0.15) {Auditoria\\{\scriptsize\texttt{auditoria}}};
\draw[linhavida] (17.20,-0.75) -- (17.20,-29.30);
\draw[msgsinc] (0.00,-1.94) -- (2.60,-1.94);
\node[rotmsg, above=1pt] at (1.30,-1.94) {verificar\\evidência};
\draw[msgsinc] (2.60,-2.98) -- (5.80,-2.98);
\node[rotmsg, above=1pt] at (4.20,-2.98) {\texttt{GET /capturas/\{id\}/verificar}\\\texttt{X-Justificativa}};
\draw[msgsinc] (5.80,-3.36) -- ++(0.35,0) -- ++(0,-0.76) -- (5.80,-4.12);
\node[rotmsg, anchor=north west, align=left] at (6.25,-3.28) {\texttt{ctx.caso}: membro (papel $\geq$ leitor);\\vale também com prazo vencido};
\draw[msgsinc] (5.80,-4.78) -- (10.00,-4.78);
\node[rotmsg, above=1pt] at (7.90,-4.78) {\texttt{verificar(captura)}};
\draw[msgsinc] (10.00,-5.16) -- ++(0.35,0) -- ++(0,-0.76) -- (10.00,-5.92);
\node[rotmsg, anchor=north west, align=left] at (10.45,-5.08) {lê \texttt{<aa>/<sha256>.gcm}\\do volume};
\draw[msgresp] (10.00,-7.51) -- (5.80,-7.51);
\node[rotmsg, above=1pt] at (7.90,-7.51) {\texttt{integra = falso},\\\texttt{arquivo\_encontrado = falso}};
\draw[msgresp] (10.00,-9.10) -- (5.80,-9.10);
\node[rotmsg, above=1pt] at (7.90,-9.10) {\texttt{integra = falso}, sem SHA-256\\recalculado};
\draw[msgsinc] (10.00,-10.03) -- ++(0.35,0) -- ++(0,-1.52) -- (10.00,-11.55);
\node[rotmsg, anchor=north west, align=left] at (10.45,-9.95) {decifra AES-256-GCM\\(AAD = SHA-256\\registrado) e\\recalcula o SHA-256};
\draw[msgresp] (10.00,-12.59) -- (5.80,-12.59);
\node[rotmsg, above=1pt] at (7.90,-12.59) {\texttt{integra} = (recalculado\\= registrado)};
\draw[frag] (2.20,-6.17) rectangle (12.50,-12.84);
\node[rotfrag, anchor=north west] (fr) at (2.20,-6.17) {alt};
\draw[frag] (fr.south west) -- ([xshift=-0.12cm]fr.south east) -- ([yshift=0.12cm]fr.south east) -- (fr.north east);
\draw[frag, dashed] (2.20,-7.76) -- (12.50,-7.76);
\draw[frag, dashed] (2.20,-9.35) -- (12.50,-9.35);
\node[rotguarda, anchor=north west] at (3.10,-6.20) {[arquivo ausente]};
\node[rotguarda, anchor=north west] at (2.32,-7.79) {[etiqueta GCM inválida (arquivo alterado ou outra chave)]};
\node[rotguarda, anchor=north west] at (2.32,-9.38) {[decifrou]};
\draw[msgsinc] (5.80,-13.60) -- (17.20,-13.60);
\node[rotmsg, above=1pt] at (11.50,-13.60) {registrar \texttt{captura.verificar}; commit; âncora};
\draw[msgresp] (5.80,-14.26) -- (2.60,-14.26);
\node[rotmsg, above=1pt] at (4.20,-14.26) {íntegra?, SHA-256 registrado e recalculado};
\draw[msgsinc] (0.00,-15.50) -- (2.60,-15.50);
\node[rotmsg, above=1pt] at (1.30,-15.50) {verificar\\a cadeia};
\draw[msgsinc] (2.60,-16.54) -- (5.80,-16.54);
\node[rotmsg, above=1pt] at (4.20,-16.54) {\texttt{GET /auditoria/verificar}\\(sem restrição de perfil)};
\draw[msgsinc] (5.80,-17.20) -- (17.20,-17.20);
\node[rotmsg, above=1pt] at (11.50,-17.20) {\texttt{verificar(sessão, chave HMAC, âncora)}};
\draw[msgsinc] (17.20,-17.86) -- (15.00,-17.86);
\node[rotmsg, above=1pt] at (16.10,-17.86) {\texttt{ler()}};
\draw[msgresp] (15.00,-18.52) -- (17.20,-18.52);
\node[rotmsg, above=1pt] at (16.10,-18.52) {\texttt{seq}, \texttt{hash}, MAC};
\draw[msgsinc] (17.20,-18.90) -- ++(0.35,0) -- ++(0,-0.76) -- (17.20,-19.66);
\node[rotmsg, anchor=north west, align=left] at (17.65,-18.82) {confere o MAC\\da âncora};
\draw[msgsinc] (17.20,-20.87) -- (12.80,-20.87);
\node[rotmsg, above=1pt] at (15.00,-20.87) {\texttt{SELECT} das entradas seguintes};
\draw[msgresp] (12.80,-21.53) -- (17.20,-21.53);
\node[rotmsg, above=1pt] at (15.00,-21.53) {lote};
\draw[msgsinc] (17.20,-21.91) -- ++(0.35,0) -- ++(0,-3.04) -- (17.20,-24.95);
\node[rotmsg, anchor=north west, align=left] at (17.65,-21.83) {para cada entrada:\\\texttt{seq} contínua; versão 2;\\\texttt{hash\_anterior} = hash\\da entrada anterior;\\HMAC-SHA256(chave,\\JSON canônico) = \texttt{hash};\\na \texttt{seq} da âncora,\\\texttt{hash} = hash ancorado};
\draw[frag] (12.35,-19.91) rectangle (21.80,-25.20);
\node[rotfrag, anchor=north west] (fr) at (12.35,-19.91) {loop};
\draw[frag] (fr.south west) -- ([xshift=-0.12cm]fr.south east) -- ([yshift=0.12cm]fr.south east) -- (fr.north east);
\node[rotguarda, anchor=north west] at (13.40,-19.94) {[lotes de entradas, em ordem de \texttt{seq}, até a primeira quebra]};
\draw[msgsinc] (17.20,-25.68) -- ++(0.35,0) -- ++(0,-1.14) -- (17.20,-26.82);
\node[rotmsg, anchor=north west, align=left] at (17.65,-25.60) {âncora ausente?\\cadeia truncada\\(âncora além do fim)?};
\draw[msgresp] (17.20,-27.48) -- (5.80,-27.48);
\node[rotmsg, above=1pt] at (11.50,-27.48) {\texttt{integra}, entradas, \texttt{primeira\_quebra}, motivo};
\draw[msgsinc] (5.80,-28.14) -- (17.20,-28.14);
\node[rotmsg, above=1pt] at (11.50,-28.14) {registrar \texttt{auditoria.verificar}; commit; âncora avança};
\draw[msgresp] (5.80,-28.80) -- (2.60,-28.80);
\node[rotmsg, above=1pt] at (4.20,-28.80) {resultado da verificação};
\end{tikzpicture}}
\fonte{Elaborada pelo autor a partir de \texttt{buscador/api/rotas\_capturas.py} (\texttt{verificar\_captura}), \texttt{buscador/custodia/servico.py} (\texttt{Custodia.verificar}), \texttt{buscador/api/rotas\_auditoria.py} (\texttt{verificar\_auditoria}), \texttt{buscador/auditoria/servico.py} (\texttt{verificar}, \texttt{Ancora}) e, no front-end, \texttt{src/componentes/CartaoAlerta.tsx} e \texttt{src/telas/abas/AbaAuditoria.tsx}.}
\end{figure}

% Diagrama de máquina de estados (UML 2) do alerta. Elementos e fontes no código:
% docs/sistema/diagramas-fontes.md, seção 7.
\begin{figure}[htbp]
\centering
\caption{Diagrama de máquina de estados do alerta}
\label{fig:uml-estado-alerta}
\resizebox{\textwidth}{!}{%
\begin{tikzpicture}[
  font=\footnotesize,
  estado/.style={draw, rounded corners=6pt, minimum width=2.7cm, minimum height=0.9cm,
    align=center, fill=white},
  inicial/.style={circle, fill=black, inner sep=0pt, minimum size=7pt},
  final/.style={circle, draw, inner sep=0pt, minimum size=11pt,
    path picture={\fill (path picture bounding box.center) circle (3pt);}},
  trans/.style={-{Straight Barb[length=6pt]}},
  rot/.style={font=\scriptsize, align=center, inner sep=1.5pt, fill=white},
  nota/.style={align=left, font=\scriptsize, inner sep=4pt}
]
\newcommand{\dobraest}[1]{%
  \draw (#1.north west) -- ([xshift=-6pt]#1.north east) -- ([yshift=-6pt]#1.north east)
        -- (#1.south east) -- (#1.south west) -- cycle;
  \draw ([xshift=-6pt]#1.north east) |- ([yshift=-6pt]#1.north east);}
% ----- pendente -----
\node[inicial] (i0) at (-0.6,0.3) {};
\node[estado] (pend) at (-0.6,-2.6) {\textbf{Pendente}\\{\scriptsize\texttt{situacao = pendente}}\\{\scriptsize\texttt{ativo = verdadeiro}}};
\node[final] (fim) at (-0.6,-6.4) {};
\draw[trans] (i0) -- node[rot, right=2pt, align=left]{tipologia detecta\\(chave nova) / cria} (pend);
\draw[trans] (pend) -- node[rot, right=2pt, align=left]{reexecução não detecta\\{[}sem decisão{]} / apaga o alerta} (fim);
\draw[trans] (pend.west) .. controls +(-1.8,0.9) and +(-1.8,-0.9) .. node[rot, left=1pt, align=right]{reexecução detecta /\\atualiza motivo,\\evidências e\\parâmetros} (pend.west);
% ----- composto: decidido -----
\node[estado] (conf) at (6.1,-0.9) {\textbf{Confirmado}\\{\scriptsize\texttt{situacao = confirmado}}};
\node[estado] (desc) at (6.1,-4.4) {\textbf{Descartado}\\{\scriptsize\texttt{situacao = descartado}}};
\draw[trans] ([xshift=-0.45cm]conf.south) -- node[rot, left=1pt]{revisão} ([xshift=-0.45cm]desc.north);
\draw[trans] ([xshift=0.45cm]desc.north) -- node[rot, right=1pt]{revisão} ([xshift=0.45cm]conf.north |- conf.south);
\node[inicial] (i1) at (9.2,-0.9) {};
\node[estado] (ativo) at (11.5,-0.9) {\textbf{Ativo}\\{\scriptsize\texttt{ativo = verdadeiro}}};
\node[estado] (inat)  at (11.5,-4.4) {\textbf{Inativo}\\{\scriptsize\texttt{ativo = falso}}};
\draw[trans] (i1) -- (ativo.west);
\draw[trans] ([xshift=-0.45cm]ativo.south) -- node[rot, left=1pt, align=right]{reexecução\\não detecta} ([xshift=-0.45cm]inat.north);
\draw[trans] ([xshift=0.45cm]inat.north) -- node[rot, right=1pt, align=left]{reexecução\\detecta} ([xshift=0.45cm]ativo.south);
\node[draw, rounded corners=8pt, inner sep=0.4cm, fit={(conf)(desc)(ativo)(inat)(i1) ($(conf.north)+(0,0.75)$) ($(inat.east)+(0.2,0)$)}] (dec) {};
\node[anchor=north, font=\footnotesize\bfseries] at (dec.north) {Decidido};
\draw ([yshift=-0.55cm]dec.north west) -- ([yshift=-0.55cm]dec.north east);
\draw[dashed] (8.75,0 |- dec.north) ++(0,-0.55) -- (8.75,0 |- dec.south);
\node[anchor=north west, font=\scriptsize\itshape] at ($(dec.north west)+(0.1,-0.6)$) {situação};
\node[anchor=north west, font=\scriptsize\itshape] at ($(8.75,0 |- dec.north)+(0.1,-0.6)$) {detecção};
% ----- entradas no composto -----
\draw[trans] (pend.east) -- node[rot, sloped, above, pos=0.45]{decidir(\texttt{confirmado},\\justificativa)} (conf.west);
\draw[trans] (pend.east) -- node[rot, sloped, below, pos=0.45]{decidir(\texttt{descartado},\\justificativa)} (desc.west);
% ----- nota -----
\node[nota, text width=12.6cm, anchor=north west] (nt) at ($(fim.south west)+(1.0,0.75)$) {%
  Toda decisão exige papel de investigador ou responsável e o caso dentro do prazo de retenção
  (\texttt{caso\_aberto}); a justificativa é obrigatória. Cada decisão, inclusive a revisão (que
  pode repetir a decisão anterior), acrescenta uma linha a \texttt{alerta\_decisoes} (só de
  acréscimo) e a entrada \texttt{alerta.decidir} à trilha de auditoria. Não há transição de volta a
  Pendente. A entrada em Decidido leva a região de detecção ao estado Ativo; um alerta decidido
  nunca é apagado.};
\dobraest{nt}
\end{tikzpicture}}
\fonte{Elaborada pelo autor a partir de \texttt{buscador/tipologias/motor.py} (\texttt{executar}),
\texttt{buscador/api/rotas\_casos.py} (\texttt{decidir\_alerta}), \texttt{buscador/api/esquemas.py}
(\texttt{DecisaoEntrada}) e \texttt{buscador/modelos.py} (\texttt{Alerta}, \texttt{AlertaDecisao}).}
\end{figure}

% Diagrama de máquina de estados (UML 2) do caso. Elementos e fontes no código:
% docs/sistema/diagramas-fontes.md, seção 8.
\begin{figure}[htbp]
\centering
\caption{Diagrama de máquina de estados do caso, com o que cada estado permite}
\label{fig:uml-estado-caso}
\resizebox{\textwidth}{!}{%
\begin{tikzpicture}[
  font=\footnotesize,
  estado/.style={draw, rounded corners=6pt, rectangle split, rectangle split parts=2,
    align=left, text width=#1, inner sep=4pt, fill=white},
  inicial/.style={circle, fill=black, inner sep=0pt, minimum size=7pt},
  trans/.style={-{Straight Barb[length=6pt]}},
  rot/.style={font=\scriptsize, align=center, inner sep=1.5pt, fill=white},
  nota/.style={align=left, font=\scriptsize, inner sep=4pt}
]
\newcommand{\dobraest}[1]{%
  \draw (#1.north west) -- ([xshift=-6pt]#1.north east) -- ([yshift=-6pt]#1.north east)
        -- (#1.south east) -- (#1.south west) -- cycle;
  \draw ([xshift=-6pt]#1.north east) |- ([yshift=-6pt]#1.north east);}
\node[estado=6.2cm, anchor=north] (ativo) at (0,0) {%
  \centering\textbf{Ativo}\par{\scriptsize hoje $\leq$ \texttt{retencao\_ate}\par ou caso sem prazo\par}
  \nodepart{second}\raggedright\scriptsize
  \textbf{Permite}, além do que Somente leitura permite:\\
  -- coletar e recoletar;\\
  -- buscar entidades, ver detalhe, grafo e caminho;\\
  -- exportar o grafo (GraphML, CSV);\\
  -- executar tipologias e decidir alertas;\\
  -- ler o texto de captura, extrair com IA e decidir sugestões;\\
  -- anotar; criar e alterar tarefas;\\
  -- criar, reativar e executar monitoramentos;\\
  -- o agendador executa os monitoramentos ativos.};
\node[estado=7.4cm, anchor=north] (leit) at (10.6,0) {%
  \centering\textbf{Somente leitura}\par{\scriptsize hoje $>$ \texttt{retencao\_ate}\par (prazo de retenção vencido)\par}
  \nodepart{second}\raggedright\scriptsize
  \textbf{Permite} (sem \texttt{caso\_aberto}):\\
  -- ver o caso; listar coletas, comparações, alertas, membros,
  anotações, tarefas, avisos, monitoramentos e sugestões;\\
  -- ver, comparar e verificar capturas; gerar o relatório;\\
  -- gerenciar membros (papel responsável);\\
  -- desativar monitoramento; marcar aviso como lido.\\[3pt]
  \textbf{Recusa} com 403 e \texttt{caso.retencao\_recusa} na trilha:\\
  -- tudo o que só Ativo permite.\\[3pt]
  \textbf{Agendador}: execução \texttt{recusada\_retencao}, monitoramento
  desativado e \texttt{monitoramento.retencao\_recusa} na trilha.};
\node[inicial] (i0) at ($(ativo.north)+(0,1.3)$) {};
\draw[trans] (i0) -- node[rot, right=2pt, align=left]{criar caso (\texttt{finalidade}, \texttt{retencao\_ate}) /\\
  criador vira responsável; \texttt{caso.criar} na trilha} (ativo.north);
\draw[trans] (ativo.east |- leit.text split) -- node[rot, above, align=center]{when(hoje UTC $>$\\\texttt{retencao\_ate})}
  (leit.west |- leit.text split);
\node[nota, text width=16.0cm, anchor=north west] (nt) at ($(ativo.south west |- leit.south)+(0,-0.35)$) {%
  Nos dois estados vale também o papel no caso: leitura exige leitor; alteração, investigador;
  membros, responsável. O prazo (\texttt{retencao\_ate}) só é gravado na criação e o último dia
  ainda vale; por isso não há volta a Ativo. Não há estado de exclusão: nada é apagado pelo
  sistema, e o descarte depois do prazo fica a cargo do órgão responsável.};
\dobraest{nt}
\end{tikzpicture}}
\fonte{Elaborada pelo autor a partir de \texttt{buscador/retencao.py}, \texttt{buscador/api/deps.py}
(\texttt{Contexto.caso\_aberto}), das rotas em \texttt{buscador/api/rotas\_*.py} (quais chamam
\texttt{caso\_aberto} e quais só \texttt{Contexto.caso}) e de \texttt{buscador/api/monitoramento.py}.}
\end{figure}


\textbf{Tipologias.} São regras explícitas, sem aprendizado de máquina,
cada uma com severidade e política de erro declaradas (decisão do autor, a
validar com o orientador): dez sobre contratações (T1 a T10), onze de
lavagem de dinheiro e uma de combinação
(Seção~\ref{sec:desenv-lavagem}). Todo alerta traz a regra, o motivo em texto e as
evidências ligadas às capturas. As tipologias que comparam pessoas (T7, sócio
servidor; T9, doador contratado) nunca afirmam identidade, porque o CPF dos
sócios nos dados abertos do CNPJ é descaracterizado: geram ``possível
coincidência'', para o investigador conferir.

% ===========================================================================
% Fonte desta seção: código do back-end na develop (buscador/tipologias/lavagem.py,
% combinacao.py, politicas.py e legislacao.py; buscador/importacao_judicial/;
% buscador/api/rotas_importacao_judicial.py) e docs/mvp-diario.md (17/11/2026).
\section{Tipologias de lavagem e importação sob ordem judicial}
\label{sec:desenv-lavagem}

As tipologias de lavagem (códigos L) seguem o mesmo contrato das T: cada
uma é uma função sobre os dados do caso, e o alerta traz a regra, o motivo
em texto, os números observados e as evidências ligadas às capturas. As T
sinalizam possíveis infrações antecedentes em contratações; as L sinalizam
possível ocultação pela interposição de pessoas e empresas
(Seção~\ref{sec:intro-contexto}). Todas são indício cadastral, e o alerta
diz isso. Nas regras que cruzam pessoas vale a regra da T7: os 6 dígitos
visíveis do CPF e o nome normalizado iguais dão ``pessoa provável'', nunca
identificação, e os registros não são fundidos. Nenhuma regra usa
sobrenome. Os limiares são decisão do autor e ficam na configuração, salvo
os limites de receita da L6, que vêm de lei. A severidade e a política de
erro de cada tipologia também são decisão do autor. A
Tabela~\ref{tab:tipologias-lavagem} resume as regras com os valores padrão.

\begin{small}
\begin{longtable}{p{1.3cm}p{8.4cm}p{4.6cm}}
\caption{Tipologias de lavagem e alerta de combinação implementados}\label{tab:tipologias-lavagem}\\
\hline
\textbf{Código} & \textbf{Regra (valores padrão)} & \textbf{Base} \\ \hline\endhead
L1 & Mesmo endereço normalizado, com o complemento, em 5 ou mais empresas com estabelecimento ativo, com ou sem contrato; o alerta diz se 3 delas iniciaram a atividade em até 180 dias & GAFI/Egmont, anexo E, itens 4 e 11; Velasco, seç.~3.1 \\
L3 & Mesma pessoa provável na administração (como sócia ou representante legal) de 5 ou mais empresas do caso, 3 delas abertas em até 12 meses & GAFI/Egmont, anexo E, itens 4 e 10; BCB, art.~1º, III, h \\
L4 & Empresa só com sócios pessoa jurídica no primeiro nível do quadro e com cadeia de 2 ou mais níveis de sócios pessoa jurídica; o alerta diz onde a cadeia termina (pessoa física, sócio estrangeiro ou empresa sem quadro nos dados) & GAFI/Egmont, anexo E, item 8; GAFI, §~47 \\
L5 & Nível 1: empresa do caso com sócio estrangeiro. O país não é avaliado & GAFI/Egmont, anexo E, item 4 \\
L6 & Soma dos contratos assinados no ano por MEI, microempresa ou empresa de pequeno porte acima do limite de receita do regime naquele ano & LC 123, arts.~3º e 18-A; LC 139, art.~7º; LC 155, art.~11; LC 238, arts.~1º, 2º e 10; Velasco, seç.~3.2 \\
L7 & Parte (a): nenhum sócio atual entrou até 30 dias depois do início da atividade, e há contrato até 12 meses depois da última entrada & GAFI/Egmont, anexo E, item 7 \\
L8 & Contratada que dispara a T1 (aberta pouco antes do contrato) e a T7 (sócio que possivelmente é servidor) & COAF, tipologia 2.7 \\
L10 & Contratada sem sanção própria que compartilha sócio-administrador provável ou endereço com empresa no CEIS ou no CNEP (sanção vigente na assinatura) ou no CEPIM; ou sócio de contratada ou de convenente que possivelmente coincide com pessoa do CEAF & GAFI/Egmont, anexo E, item 3; Lei 12.846, art.~5º, IV, e \\
L14 & Sinal fraco: todos os sócios-administradores com até 20 anos ou com mais de 80, e contratos acima de R\$~100 mil & GAFI/Egmont, anexo E, itens 3 e 10 \\
L15 & Sinal fraco: 20 ou mais atividades econômicas (CNAE) declaradas, ou contratos em 4 ou mais categorias de objeto & Velasco, seç.~3.2 \\
L16 & Entidade sem fins lucrativos convenente com ao menos um destes sinais: aberta menos de 180 dias antes da vigência; registro no CEPIM; dirigente que possivelmente é servidor do órgão concedente; dirigente que é sócio de empresa no mesmo endereço, do mesmo ramo ou recém-aberta & COAF, tipologias 2.1, 2.8 e 2.9; Velasco, seç.~5 \\
COMB & Mesma empresa ou pessoa com 2 ou mais tipologias distintas e soma de pesos de ao menos 2,0. Peso 1 por tipologia, 0,5 para sinal fraco e 0 para a L8, que é composição; a L1 e a T2 contam uma vez & Velasco, seç.~3.3 \\
\hline
\multicolumn{3}{p{14.3cm}}{\footnotesize Fonte: elaborada pelo autor, a partir do código do \textit{back-end}. Limiares, pesos, severidade e política de erro: decisão do autor, configuráveis, salvo os limites legais da L6. Bases: GAFI/Egmont \cite{fatfegmont2018}; GAFI \cite{fatf2011corruption}; Velasco \cite{velasco2021}; BCB \cite{bcb2020cc4001}; COAF \cite{coaf2021casos}; Lei 12.846 \cite{brasil2013lei12846}; LC 123 \cite{brasil2006lcp123}; LC 238 \cite{brasil2026lcp238}.}\\
\end{longtable}
\end{small}

\textbf{Limites legais da L6.} A regra segue a lei, sem tolerância própria.
O MEI tem limite de R\$~81.000,00 por ano-calendário pela redação atual
\cite[art.~18-A, § 1º]{brasil2006lcp123}, de R\$~110.000,00 em 2027
\cite[art.~2º]{brasil2026lcp238} e de R\$~140.000,00 a partir de 2028
\cite[arts.~1º e 10, I]{brasil2026lcp238}. A microempresa vai até
R\$~360.000,00 e a empresa de pequeno porte até R\$~4.800.000,00
\cite[art.~3º, I e II]{brasil2006lcp123}. Esses valores de MEI e de empresa
de pequeno porte valem desde 2018 \cite[art.~11, III]{brasil2016lcp155}; de
2012 a 2017 valiam R\$~60.000,00 para o MEI e R\$~3.600.000,00 para a
empresa de pequeno porte \cite[art.~7º]{brasil2011lcp139}. Anos anteriores a
2012 não são avaliados, e o alerta diz isso.

No ano de início de atividade, o limite é proporcional. Para o MEI, é o
valor mensal da lei vezes os meses entre o início e o fim do ano, com a
fração de mês contada como mês inteiro \cite[art.~18-A, § 2º]{brasil2006lcp123};
a partir de 2027, os valores mensais mudam \cite[arts.~1º e 2º]{brasil2026lcp238}.
Para a microempresa e a empresa de pequeno porte, o limite é proporcional
aos meses de atividade, inclusive as frações \cite[art.~3º, § 2º]{brasil2006lcp123}.

O alerta dispara com qualquer valor acima do limite; igual ao limite não
dispara. Ele informa se o excesso fica até 20\% ou acima disso, porque a lei
usa essa faixa para definir quando o desenquadramento produz efeitos
\cite[art.~3º, §§ 9º-A e 12; art.~18-A, § 7º, III e IV]{brasil2006lcp123}.
Duas escolhas são decisão do autor:
\begin{itemize}
	\item a lei mede a receita bruta, que não é dado aberto; a regra usa a
	soma dos contratos públicos assinados no ano (valor inicial, ou o global
	quando não há outro) como indicador mínimo dessa receita;
	\item o regime do MEI é o da tabela Simples do CNPJ na data de cada
	assinatura, mas o porte da microempresa e da empresa de pequeno porte é
	o do cadastro atual.
\end{itemize}

\textbf{O que não entrou.} A L2 (contato compartilhado), o nível~2 da L5
(lista de jurisdições, que depende de lista oficial ainda não conferida) e
a parte (b) da L7 (troca de sócios entre duas coletas) não foram
implementados. A L8 e a L16 não usam eleito nem pessoa exposta
politicamente, porque isso depende da L9. A L9, a L11, a L12 e a L13
aguardam decisão do orientador (Seção~\ref{sec:desenv-limitacoes}).

\textbf{Combinação.} A combinação lê os alertas do caso, e não as
entidades. Só entram alertas ativos e não descartados pelo investigador, e
cada tipologia conta uma vez por entidade. O alerta lista os alertas que o
compõem, com as evidências, e diz que não é escore de culpa, e sim uma pista
para priorizar a análise. A base é a observação de
\citeonline[seç.~3.3]{velasco2021} de que os padrões, isolados, não
representam risco significativo, e de que o risco cresce quando eles
aparecem combinados.

\textbf{Importação sob ordem judicial (RF10).} É um módulo separado, desligado
por padrão. Com ele desligado, só a rota de estado responde. Com ele
ligado, o fluxo é este:
\begin{enumerate}
	\item um usuário de perfil autorizado registra a decisão judicial antes
	de qualquer arquivo: processo, decisão, juízo, data, período, contas e
	escopo;
	\item cada arquivo é validado pelo leiaute do Anexo I da Instrução
	Normativa BCB nº~636/2025, itens I a V (agências, contas, titulares,
	extrato e origem-destino), em texto separado por tabulação, sem
	cabeçalho \cite[arts.~2º e 3º]{bcb2025in636}. A linha fora do leiaute é
	rejeitada, e o relatório indica a linha, o campo e o motivo, sem o
	conteúdo;
	\item o que estiver fora das contas ou do período autorizados fica em
	quarentena, sinalizado;
	\item os registros vão para uma base própria, cifrados com chave própria
	(AES-256-GCM, como as capturas), com a marca da decisão judicial. Nada
	disso vira entidade, aresta, alerta, relatório ou exportação;
	\item cada arquivo recebido entra na custódia, com SHA-256, tamanho,
	origem e carimbo de tempo;
	\item toda chamada às rotas do módulo dentro de um caso, inclusive as de
	leitura, entra na
	trilha de auditoria com a justificativa, sem número de conta, documento
	nem conteúdo de registro.
\end{enumerate}
O módulo só aceita pastas com um arquivo-marcador de dado sintético. Dado
real nunca foi usado. Ficaram de fora o RIF, o Anexo II da instrução
normativa (os códigos de tipo de lançamento, validados só como número), o
leiaute anterior a ela, a camada do grafo, os alertas financeiros e o
descarte ao fim da retenção. Que o SIMBA entregue os dados nesse leiaute é
inferência do autor, pelos nomes dos arquivos, e não foi conferido.

\textbf{\textit{Front-end}.} O catálogo de tipologias passou a incluir as L
e a combinação. Cada cartão de alerta traz o aviso ``Pista, não
conclusão'', e o caso ganhou uma aba de ordem judicial para o RF10. A tela
nova não foi conferida em navegador nem contra a API real.

% ===========================================================================
\section{Interface de programação}
\label{sec:desenv-api}

A Tabela~\ref{tab:api} lista os 59 \textit{endpoints} da API, extraídos da
especificação OpenAPI gerada pelo próprio FastAPI. Toda rota de caso exige
sessão (ou token de desenvolvimento, só no modo de desenvolvimento),
justificativa no cabeçalho, ser membro do caso com o papel adequado e caso
dentro do prazo de retenção para escrever. Cada chamada entra na trilha de
auditoria.

% Gerado por docs/sistema/gera_api_tex.py a partir do OpenAPI do back-end. Não editar à mão.
\begin{small}
\begin{longtable}{lp{6.4cm}p{6.6cm}}
\caption{Endpoints da API do protótipo}\rotulounico{tab:api}\\
\hline
\textbf{Método} & \textbf{Caminho} & \textbf{Descrição} \\ \hline\endhead
GET & \texttt{/casos/\{caso\_id\}/alertas} & Alertas do caso por severidade (alta primeiro), com a política de erro da tipologia (RF14) e a situação de cada um (RF12). \\
POST & \texttt{/casos/\{caso\_id\}/alertas/executar} & Executa as tipologias T1 a T10 e reconcilia os alertas do caso (RF04). \\
POST & \texttt{/casos/\{caso\_id\}/alertas/\{alerta\_id\}/decisao} & Confirma ou descarta o alerta, com justificativa (RF12, H12). \\
GET & \texttt{/casos/\{caso\_id\}/anotacoes} & Anotações do caso, da mais recente para a mais antiga; com ``alvo\_tipo`` e ``alvo\_id``, só as de uma entidade ou de um alerta. \\
POST & \texttt{/casos/\{caso\_id\}/anotacoes} & Anota uma entidade ou um alerta do caso. \\
GET & \texttt{/casos/\{caso\_id\}/avisos} & Avisos de novidade do caso, do mais recente para o mais antigo. \\
POST & \texttt{/casos/\{caso\_id\}/avisos/\{aviso\_id\}/lido} & Marca um aviso de novidade como lido (H18). \\
GET & \texttt{/casos/\{caso\_id\}/caminho} & Caminho entre duas entidades, por busca bidirecional no banco (RF03, H29). \\
POST & \texttt{/casos/\{caso\_id\}/capturas/\{captura\_id\}/extracao-ia} & Roda o modelo local sobre o texto da captura. \\
GET & \texttt{/casos/\{caso\_id\}/capturas/\{captura\_id\}/texto} & O texto do documento capturado, para conferir o trecho citado ao lado do original. \\
GET & \texttt{/casos} & Só os casos de que o usuário é membro (H23). \\
POST & \texttt{/casos} & Cria um caso com finalidade e prazo de retenção obrigatórios; o criador vira responsável (H17, H23). \\
GET & \texttt{/casos/\{caso\_id\}} & Detalhe do caso, com o papel do usuário e as contagens. \\
GET & \texttt{/casos/\{caso\_id\}/coletas} & Histórico das coletas do caso, da mais recente para a mais antiga (RNF12, H15). \\
POST & \texttt{/casos/\{caso\_id\}/coletas} & Executa uma coleta passiva numa fonte, grava as capturas na custódia e resolve as entidades (RF01, RF07). \\
GET & \texttt{/casos/\{caso\_id\}/coletas/\{coleta\_id\}/comparacao} & Compara as capturas de uma coleta com as de outra do mesmo caso e da mesma fonte. \\
POST & \texttt{/casos/\{caso\_id\}/coletas/\{coleta\_id\}/recoletar} & Repete uma coleta do caso (mesma fonte e mesmos parâmetros) e compara as capturas novas com as da coleta original (RF07, H14). \\
GET & \texttt{/casos/\{caso\_id\}/entidades} & Busca entidades do caso por nome, identificador e tipo (RF11). \\
GET & \texttt{/casos/\{caso\_id\}/entidades/\{entidade\_id\}} & Detalhe da entidade, com fontes e vínculos. \\
GET & \texttt{/casos/\{caso\_id\}/exportacao} & Grafo do caso em GraphML, CSV de entidades, CSV de vínculos ou ZIP com os três. \\
GET & \texttt{/casos/\{caso\_id\}/extracoes-ia} & Execuções da extração no caso (inclusive as que falharam), da mais recente para a mais antiga, com modelo, versão, parâmetros e os itens recusados de cada uma. \\
GET & \texttt{/casos/\{caso\_id\}/grafo} & Vizinhança de uma entidade em N níveis, consultada no banco (RF03, H29). \\
GET & \texttt{/casos/\{caso\_id\}/membros} & Membros do caso e seus papéis (H23). \\
POST & \texttt{/casos/\{caso\_id\}/membros} & Acrescenta um membro ao caso ou muda o papel de quem já é membro. \\
POST & \texttt{/casos/\{caso\_id\}/membros/\{usuario\}/remover} & Tira o membro do caso (ele deixa de ver o caso na hora). \\
GET & \texttt{/casos/\{caso\_id\}/monitoramentos} & Entidades monitoradas no caso (H18). \\
POST & \texttt{/casos/\{caso\_id\}/monitoramentos} & Marca uma entidade do caso para acompanhamento. \\
POST & \texttt{/casos/\{caso\_id\}/monitoramentos/\{monitoramento\_id\}/executar} & Executa o monitoramento agora (o agendador faz o mesmo na periodicidade escolhida). \\
POST & \texttt{/casos/\{caso\_id\}/monitoramentos/\{monitoramento\_id\}/situacao} & Ativa ou desativa o monitoramento e, opcionalmente, muda a periodicidade. \\
GET & \texttt{/casos/\{caso\_id\}/relatorio} & Relatório do caso em JSON ou Markdown: entidades, alertas e decisões, coletas, custódia (RF09). \\
GET & \texttt{/casos/\{caso\_id\}/sugestoes-ia} & Sugestões da extração por IA, pendentes ou decididas (H22). \\
POST & \texttt{/casos/\{caso\_id\}/sugestoes-ia/\{sugestao\_id\}/decisao} & O investigador aceita ou recusa a sugestão, com justificativa. \\
GET & \texttt{/casos/\{caso\_id\}/tarefas} & Tarefas do caso (H37). \\
POST & \texttt{/casos/\{caso\_id\}/tarefas} & Cria uma tarefa do caso, com histórico (H37). \\
GET & \texttt{/casos/\{caso\_id\}/tarefas/\{tarefa\_id\}} & Detalhe da tarefa, com o histórico de alterações (H37). \\
POST & \texttt{/casos/\{caso\_id\}/tarefas/\{tarefa\_id\}} & Altera os campos enviados (título, descrição, responsável, situação, prazo). \\
GET & \texttt{/auditoria} & Trilha de auditoria (delegado e administrador) (RNF02). \\
GET & \texttt{/auditoria/verificar} & Verifica o encadeamento por HMAC da trilha e a âncora (RNF02). \\
POST & \texttt{/auth/login} & Primeira etapa: usuário e senha. \\
POST & \texttt{/auth/logout} & Encerra a sessão atual na hora. \\
POST & \texttt{/auth/mfa} & Segunda etapa: o código TOTP, com o token pendente em ``Authorization``. \\
POST & \texttt{/auth/senha} & Troca a própria senha (exige a atual) e encerra as outras sessões do usuário. \\
GET & \texttt{/auth/sessao} & Situação da sessão atual (etapa, expiração). \\
GET & \texttt{/capturas/\{captura\_id\}} & Metadados de custódia de uma captura: hash, origem, instante, conector e versão (RF07). \\
GET & \texttt{/capturas/\{captura\_id\}/comparar} & Compara a captura com outra da mesma origem (RF07, H14). \\
GET & \texttt{/capturas/\{captura\_id\}/verificar} & Recalcula o SHA-256 da captura (decifrando, se cifrada) e compara com o registrado (RF07). \\
GET & \texttt{/health} & Verificação de saúde do serviço, usada pelo Docker Compose. \\
GET & \texttt{/ia} & Diz se a extração por IA está ligada e com qual provedor e modelo (não é dado de caso). \\
GET & \texttt{/me} & Dados do usuário autenticado (login e perfil). \\
GET & \texttt{/tipologias} & Catálogo das tipologias (T1 a T10): regra, base, política de erro e a configuração em vigor. \\
GET & \texttt{/usuarios} & Lista as contas de usuário (só administrador). \\
POST & \texttt{/usuarios} & Cria a conta com uma senha inicial. \\
POST & \texttt{/usuarios/\{usuario\}} & Ativa ou desativa, muda o perfil, redefine a senha ou apaga o segundo fator (o usuário cadastra outro no próximo login). \\
\hline
\multicolumn{3}{p{14cm}}{\footnotesize Fonte: gerada pelo autor a partir do OpenAPI do back-end (FastAPI). A descrição é a primeira frase da documentação de cada rota no código; para as rotas sem documentação, foi escrita pelo autor a partir do código (\texttt{docs/sistema/api-descricoes.json}).}\\
\end{longtable}
\end{small}


% ===========================================================================
\section{Decisões de arquitetura}
\label{sec:desenv-decisoes}

A Tabela~\ref{tab:adr} registra as principais decisões de arquitetura, com o
contexto, a decisão, as alternativas consideradas e a base de cada uma.
Quando a base é uma lei, um requisito ou um achado das revisões, ela é citada;
quando não há fonte externa, a decisão é do autor.

\begin{small}
\begin{longtable}{p{2.6cm}p{4.4cm}p{3.6cm}p{3.6cm}}
\caption{Registro de decisões de arquitetura}\rotulounico{tab:adr}\\
\hline
\textbf{Decisão} & \textbf{O que foi decidido} & \textbf{Alternativas} & \textbf{Base} \\ \hline\endhead
Coleta só passiva & Ler só dados abertos, sem interagir com alvos nem usar contas & Coleta ativa, perfis controlados & Lei 12.850/2013, art.~10-A \cite{brasil2013lei12850}; histórias H11 e H26 \\
Homônimos & Sugestão de unificação, nunca fusão automática, quando só o nome coincide & Fusão por nome & CPF descaracterizado no CNPJ (Receita Federal); RSL~1, similaridade é pista (SB1-0781) \\
Cadeia de custódia & Cada carga guardada como veio, com SHA-256, instante UTC, origem, conector e versão; verificação sob demanda & Guardar só os dados extraídos & RF07; RSL~1, QP5 (A3, TA4) \\
Trilha de auditoria & Só de acréscimo, encadeada por HMAC com chave fora do banco, com gatilhos e âncora & Registro comum em tabela & RNF02; RSL~1 (A4, TD6) \\
Papéis do banco & Papel dono (migração) separado do papel da aplicação, que só tem os privilégios necessários & Um só usuário de banco & Menor privilégio; decisão do autor \\
Alerta como pista & Alerta exige confirmação ou descarte com justificativa; histórico guardado & Alerta como conclusão & RF12; RSL~1 (A2, TA3; A4, TA3) \\
Regras em vez de IA nas tipologias & Vinte e duas tipologias por regras explícitas (dez de contratações, onze de lavagem e a combinação), com motivo e evidências & Classificador treinado & RNF05; nenhum estudo da RSL~1 trata de fraude em contratações nem do Brasil (Seção~\ref{sec:requisitos-mercado}) \\
IA só local e opcional & Extração por modelo local, desligada por padrão; trecho literal obrigatório; decisão humana & Serviço de IA externo & Soberania dos dados; RSL~1 (A2, TA4); RSL~2 (B3, TA3) \\
Sessão & Sessão no servidor com token opaco, expiração por duração e inatividade & JWT & Logout e desativação valem na hora; decisão do autor \\
Acesso por caso & Membros com papéis por caso; perfil global não abre caso & Acesso só por perfil & RNF01; minimização (RNF04) \\
Criptografia em repouso & AES-256-GCM das capturas, chave fora do banco, SHA-256 sobre o conteúdo original & Cifrar o disco inteiro & RNF01; decisão do autor \\
Retenção & Depois do prazo, o caso fica só leitura; nada é apagado & Apagar ao fim do prazo & RNF04; custódia e auditoria são só de acréscimo; descarte é decisão do órgão \\
Grafo no banco & Busca nível a nível no PostgreSQL, sem novo serviço & NetworkX em memória; banco de grafos & Medição de desempenho (Seção~\ref{sec:desenv-verificacao}) \\
Dados de teste & Só dados sintéticos, CPFs com dígito inválido ou gerados na execução & Amostras reais & LGPD; regra do projeto \\
Editais em PDF & Custódia sobre os bytes do PDF; texto derivado, com o SHA-256 do texto; extração em processo separado, com limites; sem OCR & Guardar só o texto extraído; OCR & Custódia (RF07); decisão do autor \\
CNPJ por recorte & \textit{Script} lê os ZIPs oficiais em fluxo e grava só os CNPJs do caso, com manifesto dos SHA-256 & Carga da base completa pelo conector & Mais de 26~GB descompactados (estimativa por proporção, 17/11/2026); decisão do autor \\
\hline
\multicolumn{4}{p{14.6cm}}{\footnotesize Fonte: elaborada pelo autor, a partir do código, dos requisitos e da síntese das revisões.}\\
\end{longtable}
\end{small}
% PENDENTE: o orientador valida as decisões marcadas como do autor.

% ===========================================================================
\section{Modelo de ameaças}
\label{sec:desenv-ameacas}

A síntese da RSL~1 mostrou que o modelo de ameaças de uma ferramenta OSINT
policial tem três alvos: o investigador, a ferramenta e os dados (A2, TA4,
Seção~\ref{sec:resultados-rsl1}). A Tabela~\ref{tab:ameacas} relaciona cada
ameaça tratada no protótipo ao controle implementado e ao teste automatizado
que o verifica.

\begin{small}
\begin{longtable}{@{}p{3.3cm}p{2.1cm}p{4.4cm}p{4.5cm}@{}}
\caption{Ameaças, controles e testes}\rotulounico{tab:ameacas}\\
\hline
\textbf{Ameaça} & \textbf{Alvo} & \textbf{Controle} & \textbf{Teste (back-end)} \\ \hline\endhead
Acesso a caso alheio & Dados & Acesso por caso com papéis; negação auditada & \nometeste{test_h23_nao_membro_nao_ve_nem_age_no_caso_e_a_negacao_e_auditada} \\
Roubo de senha & Ferramenta & Argon2id, MFA por TOTP, bloqueio após 5 falhas, código de uso único & \nometeste{test_h23_senha_errada_usuario_inexistente_e_bloqueio_por_tentativas}; \nometeste{test_h23_mfa_o_mesmo_codigo_nao_vale_duas_vezes} \\
Adulteração da prova & Dados & SHA-256 na captura e verificação; cifra autenticada & \nometeste{test_h23_adulteracao_troca_de_arquivo_e_chave_errada_sao_detectadas} \\
Adulteração da trilha & Ferramenta & HMAC encadeado, gatilhos, papel sem UPDATE e DELETE & \nometeste{test_adulteracao_via_sql_e_detectada}; \nometeste{test_gatilho_do_banco_bloqueia_update_e_delete} \\
Vazamento em repouso & Dados & AES-256-GCM das capturas, chave fora do banco & \nometeste{test_h23_captura_fica_cifrada_no_disco_e_o_sha256_e_do_conteudo_original} \\
Envio de dados a terceiros pela IA & Dados & IA desligada por padrão; só servidor local; modelo remoto recusado & \nometeste{test_h22_configuracao_so_aceita_servidor_local}; \nometeste{test_h22_provedor_ollama_recusa_modelo_remoto_ausente_e_respostas_ruins} \\
Conteúdo inventado pela IA & Ferramenta & Trecho literal, valor dentro do trecho e esquema obrigatórios & \nometeste{test_h22_trecho_inventado_valor_inventado_e_esquema_invalido_sao_recusados} \\
Exposição de CPF & Dados & Só os seis dígitos centrais em entidades, alertas, avisos e relatório; máscara na interface & \nometeste{test_rnf04_nenhum_cpf_completo_chega_a_entidades_alertas_avisos_ou_relatorio} \\
CPF completo em dado bruto & Dados & O contrato do PNCP guarda só o CPF descaracterizado do fornecedor; as capturas ficam cifradas no disco. Risco em tratamento: o CPF completo continua nas capturas brutas e na comparação entre capturas & \nometeste{test_cpf_de_fornecedor_pessoa_fisica_nao_fica_completo_no_contrato} \\
PDF malicioso & Ferramenta & Extração em processo separado, sem as variáveis de ambiente da API e encerrada no prazo; só o texto das páginas, sem JavaScript nem anexos; limites de tamanho, páginas, caracteres e tempo & \nometeste{test_javascript_e_anexo_do_pdf_nao_sao_executados_nem_lidos}; \nometeste{test_limites_de_paginas_tamanho_caracteres_e_tempo} \\
Texto do PDF diferente do da coleta & Dados & SHA-256 do texto extraído na coleta, comparado a cada nova extração & \nometeste{test_texto_extraido_diferente_do_da_coleta_e_sinalizado} \\
Fórmula maliciosa em planilha exportada & Investigador & Neutralização de fórmulas no CSV; escape de XML no GraphML & \nometeste{test_h38_csv_neutraliza_formulas_e_graphml_escapa_xml} \\
Alertar o alvo ou deixar rastro & Investigador & Só coleta passiva de dados abertos & Desenho dos conectores (só GET); anonimato operacional fora desta entrega (H26) \\
Consulta que derruba o serviço & Ferramenta & Tetos de nós na resposta e nos visitados & \nometeste{test_h29_tetos_e_erros_das_consultas} \\
Uso fora da finalidade & Dados & Justificativa em toda chamada; retenção com só leitura & \nometeste{test_rotas_de_caso_exigem_token_e_justificativa}; \nometeste{test_h17_caso_exige_finalidade_e_prazo_de_retencao} \\
\hline
\multicolumn{4}{p{14.6cm}}{\footnotesize Fonte: elaborada pelo autor, a partir do código e dos testes do \textit{back-end}.}\\
\end{longtable}
\end{small}
% PENDENTE: atualizar quando a correção do CPF/CNPJ alfanumérico terminar.
% Em 17/11/2026, outro agente corrige no back-end a exibição do CPF completo
% (capturas brutas e comparação entre capturas) e o CNPJ alfanumérico.

Antes do horizonte~1, o protótipo passou por uma auditoria de segurança
assistida por IA, e os 17 itens encontrados foram corrigidos, cada um com
teste de regressão (execuções 4 e 5 do diário, Seção~\ref{sec:desenv-processo}).
% PENDENTE do CI: ver a Seção sec:desenv-testes-ci.

% ===========================================================================
\section{Engenharia de testes e verificação}
\label{sec:desenv-verificacao}
% Fontes: testes dos dois repositórios (tests/, conftest.py, sinteticos.py,
% src/**/*.test.ts(x), src/test/), pyproject.toml, vite.config.ts,
% package.json, .github/workflows/ci.yml dos dois repositórios e
% docs/mvp-diario.md. Termos de teste conforme o SWEBOK v4.0, cap. 5.

Esta seção descreve como o protótipo foi verificado: a estratégia de testes,
as técnicas usadas, a ligação entre requisitos e testes, os portões da
integração contínua, o teste contra as fontes reais, a medição de desempenho
e as lacunas. Os termos seguem o capítulo de teste do SWEBOK v4.0
\cite{swebok4}. O mesmo guia lembra o limite de toda a seção: o teste mostra
a presença de defeitos, mas nunca a ausência deles \cite[cap.~5, seç.~1.2.8]{swebok4}.

% ---------------------------------------------------------------------------
\subsection{Estratégia}
\label{sec:desenv-testes-estrategia}

O SWEBOK distingue níveis de teste pelo alvo (unidade, integração, sistema e
aceitação) e tipos de teste pelo objetivo \cite[cap.~5, seç.~2]{swebok4}.
O teste de unidade verifica, isolado, um elemento testável em separado; o de
integração verifica as interações entre elementos; e o de sistema verifica o
comportamento do conjunto e é o lugar usual dos requisitos não funcionais,
como segurança e privacidade \cite[cap.~5, seç.~2.1]{swebok4}. A
Tabela~\ref{tab:niveis-teste} mostra como cada nível aparece no protótipo.
A escolha de testar cada nível com essas ferramentas é decisão do autor.

\begin{small}
\begin{longtable}{p{2.6cm}p{5.6cm}p{6.2cm}}
\caption{Níveis de teste do protótipo}\label{tab:niveis-teste}\\
\hline
\textbf{Nível} & \textbf{O que se verifica} & \textbf{Como} \\ \hline\endhead
Unidade & Funções isoladas: normalização de nomes e endereços, dígito verificador do CNPJ, descaracterização do CPF, código TOTP, cofre AES-256-GCM, extração de texto de PDF, conversão do grafo para o Cytoscape.js & \textit{pytest}; Vitest (por exemplo, \texttt{test\_normalizacao.py} e \texttt{privacidade.test.ts}) \\
Integração com o banco & Serviços sobre o banco real: custódia, trilha de auditoria, resolução, tipologias, migração do esquema; no PostgreSQL, também os gatilhos e os privilégios do papel da aplicação & SQLite em arquivo temporário em todo teste; PostgreSQL só quando as variáveis \texttt{TEST\_POSTGRES\_URL} e \texttt{TEST\_POSTGRES\_URL\_APP} existem; sem elas, os 8 testes de PostgreSQL são pulados (\texttt{pytest.mark.skipif}) \\
API & A aplicação inteira pela interface HTTP: autenticação, justificativa, acesso por caso, coleta, alertas, relatório e exportação & \texttt{TestClient} do FastAPI sobre a aplicação montada com os conectores simulados \\
\textit{Front-end} & Componentes e telas: o que é mostrado, o que é enviado à API, as recusas e os papéis & Vitest com jsdom e a Testing Library; cliente da API falso \\
\hline
\multicolumn{3}{p{14.4cm}}{\footnotesize Fonte: elaborada pelo autor, a partir de \texttt{tests/conftest.py}, dos testes e do \texttt{vite.config.ts}.}\\
\end{longtable}
\end{small}

Não há testes de aceitação automatizados nem testes ponta a ponta, com o
navegador e o \textit{back-end} reais juntos. A avaliação com usuários tem
protocolo (Capítulo~\ref{cap:avaliacao}), mas ainda não foi executada.

Os tipos de teste presentes são seis:
\begin{itemize}
	\item \textbf{Funcional:} as tipologias têm casos que disparam e, na
	maioria, casos que não disparam (por exemplo, \texttt{test\_\allowbreak{}t1\_\allowbreak{}dispara\_\allowbreak{}para\_\allowbreak{}empresa\_\allowbreak{}recente\_\allowbreak{}e\_\allowbreak{}nao\_\allowbreak{}para\_\allowbreak{}antiga}
	e \texttt{test\_\allowbreak{}h31\_\allowbreak{}t4\_\allowbreak{}nao\_\allowbreak{}dispara\_\allowbreak{}para\_\allowbreak{}ativa\_\allowbreak{}nem\_\allowbreak{}para\_\allowbreak{}empresa\_\allowbreak{}sem\_\allowbreak{}dado\_\allowbreak{}do\_\allowbreak{}cnpj}).
	\item \textbf{Segurança:} o SWEBOK a define como a verificação da
	confidencialidade, da integridade e da disponibilidade, incluindo o uso
	indevido, isto é, o teste negativo \cite[cap.~5, seç.~2.2.8]{swebok4}. No
	protótipo, são os testes da Tabela~\ref{tab:ameacas} e as regressões da
	auditoria de segurança (\texttt{test\_seguranca.py}).
	\item \textbf{Privacidade:} o guia trata o teste de privacidade como tipo
	próprio \cite[cap.~5, seç.~2.2.9]{swebok4}. Aqui, ele verifica que o CPF
	completo não sai da custódia por nenhum caminho da API
	(\texttt{test\_privacidade\_cpf.py}).
	\item \textbf{Regressão:} o novo teste depois de uma mudança, para mostrar
	que o que passava continua passando \cite[cap.~5, seç.~2.2.5]{swebok4}.
	Cada um dos 17 achados da auditoria de segurança e cada correção vinda das
	fontes reais ganhou um teste (\texttt{test\_fontes\_reais.py}).
	\item \textbf{Desempenho:} a verificação de que o sistema atende aos
	requisitos de desempenho \cite[cap.~5, seç.~2.2.7]{swebok4}. A medição do
	grafo (Seção~\ref{sec:desenv-desempenho}) é feita à parte, por
	\textit{script}, e não entra na suíte nem no \textit{pipeline}. Na suíte, os
	testes da H29 conferem só o resultado das consultas, e não o tempo.
	\item \textbf{Fumaça:} o teste rápido de que as funções centrais operam
	\cite[cap.~5, seç.~3.3.3]{swebok4}. Aqui, é a execução manual dos conectores
	contra as fontes reais (Seção~\ref{sec:desenv-fontes-reais}). O arquivo
	\texttt{test\_fontes\_reais.py} testa só o \textit{script} dessa execução,
	sem rede.
\end{itemize}

% ---------------------------------------------------------------------------
\subsection{Técnicas}
\label{sec:desenv-testes-tecnicas}

\textbf{Dublês de teste.} O SWEBOK descreve \textit{stubs} e
\textit{drivers} que simulam os módulos chamados e os que chamam
\cite[cap.~5, seç.~8.2]{swebok4}. Simular uma dependência dá foco e
velocidade ao teste, mas tira realismo, e manter o simulado compatível com o
original é difícil \cite{spadini2017mock}. O protótipo usa estes dublês:
\begin{itemize}
	\item \texttt{TransporteSimulado}, transporte do \texttt{httpx} que serve
	as páginas do PNCP a partir de arquivos locais, sem rede. Ele registra as
	chamadas recebidas e pode forçar respostas 503 nas primeiras chamadas,
	para testar a retentativa. As capturas feitas com ele ficam marcadas como
	sintéticas;
	\item \texttt{TransporteRotas}, transporte do mesmo tipo que responde por
	uma função, como o Portal da Transparência ou os termos do PNCP
	responderiam, inclusive com erro;
	\item \texttt{httpx.MockTransport}, que responde como o servidor Ollama
	do modelo local;
	\item \texttt{ProvedorFalso}, provedor de IA determinístico, sem modelo.
	Nos testes, ele devolve respostas prontas, inclusive inválidas ou com
	trecho inventado, para provar que o serviço as recusa;
	\item substituições pontuais com o \texttt{monkeypatch} do \textit{pytest},
	como a data corrente da retenção, para testar o último dia do prazo, e
	funções espiãs que registram se uma etapa foi chamada;
	\item no \textit{front-end}, um \texttt{fetch} falso (\texttt{vi.fn}) para o
	cliente HTTP e o \texttt{clienteFalso}, que parte do cliente simulado
	(\texttt{ClienteMock}) e troca só os métodos de cada teste.
\end{itemize}

\textbf{Dados sintéticos.} Os testes só usam dados inventados
(Tabela~\ref{tab:adr}), como recomenda a RSL~2 para o código feito com IA
(Seção~\ref{sec:resultados-rsl2-praticas}). Como todo número de CPF com
dígitos verificadores válidos pode ser de alguém, nenhum fica gravado nos
repositórios. A cada execução, o teste sorteia a base e calcula os
verificadores (\texttt{cpf\_sintetico} no \textit{back-end} e
\texttt{cpfSinteticoValido} no \textit{front-end}). O CNPJ alfanumérico
também é sorteado, com o dígito verificador da regra oficial. Os PDFs de
editais são gerados em memória, e as \textit{fixtures} das fontes vêm de
\textit{scripts} geradores (\texttt{scripts/gerar\_fixtures.py} e
\texttt{scripts/gerar\_fixtures\_h2.py}). Essa geração aleatória é a seleção
aleatória de entradas descrita no SWEBOK \cite[cap.~5, seç.~3.7.2]{swebok4}.
Não há semente registrada, o que tem um custo (Seção~\ref{sec:desenv-testes-lacunas}).

\textbf{Oráculos.} O oráculo decide se o comportamento observado está
correto \cite[cap.~5, seç.~1.2.7]{swebok4}. Barr et al.\ classificam os
oráculos em especificados, derivados, implícitos e humanos
\cite{barr2015oracle}. O protótipo usa dois deles:
\begin{itemize}
	\item \textbf{Especificados:} os vetores de teste da RFC~6238 para o
	código TOTP; o exemplo oficial do manual da Receita para o dígito
	verificador do CNPJ alfanumérico; os limites de 2026 do decreto de
	atualização, para a T5. Também são especificados os casos de cada
	tipologia, escritos a partir da regra, e as
	verificações de integridade: o teste da custódia recalcula o SHA-256 dos
	bytes com a biblioteca padrão e o compara com o da captura; na auditoria,
	o teste altera uma linha direto no banco, ou reescreve a cadeia com outra
	chave de HMAC, e exige que a verificação acuse a quebra.
	\item \textbf{Derivados:} os testes da H29 comparam, para todo nó e toda
	profundidade, a consulta no banco com a implementação anterior, em
	NetworkX, com o grafo inteiro em memória. É uma versão alternativa do
	programa usada como oráculo, perto do que Barr et al.\ chamam de
	pseudo-oráculo \cite{barr2015oracle}. Pela definição deles, porém, a
	versão alternativa deve ser produzida de forma independente, e aqui as
	duas foram escritas pelo mesmo agente.
\end{itemize}
No teste com as fontes reais e na execução com o modelo local
(Seção~\ref{sec:desenv-limitacoes}), a saída foi conferida pelos agentes de
IA, e não há registro de conferência humana. Por isso, esses casos não contam
como oráculo humano.

Não há testes baseados em propriedades: nenhum dos repositórios usa
bibliotecas desse tipo.

% ---------------------------------------------------------------------------
\subsection{Rastreabilidade entre requisitos e testes}
\label{sec:desenv-testes-rastreabilidade}
% Gerado por docs/sistema/gera_rastreabilidade.py. Não editar à mão.
\newcommand{\rastTestesBack}{318}
\newcommand{\rastTestesFront}{244}
\newcommand{\rastLigBack}{288}
\newcommand{\rastLigFront}{176}
\newcommand{\rastNomeBack}{111}
\newcommand{\rastNomeFront}{108}
\newcommand{\rastExcluidos}{48}
\newcommand{\rastSemVinculoBack}{30}
\newcommand{\rastSemVinculoFront}{20}
\newcommand{\rastReqSemTeste}{7}
\newcommand{\rastElegiveisFront}{196}


Uma medida de cobertura citada pelo SWEBOK é a porcentagem dos requisitos
funcionais exercitados pelos testes \cite[cap.~5, seç.~4]{swebok4}. Para
medi-la, um \textit{script} (\texttt{docs/sistema/gera\_rastreabilidade.py})
lê os testes dos dois repositórios, sem executá-los, e liga cada teste aos
requisitos do Capítulo~\ref{cap:requisitos}. A ligação tem três níveis, nesta
ordem:
\begin{enumerate}
	\item \textbf{Nome:} o identificador está no próprio teste. No
	\textit{back-end}, o nome começa por \nometeste{test_h23_},
	\nometeste{test_rf02_} ou \nometeste{test_rnf04_}; no \textit{front-end}, o
	título do teste ou do grupo cita H, RF ou RNF. A história vira requisitos
	pelo próprio \textit{backlog} (Tabela~\ref{tab:backlog}).
	\item \textbf{Padrão:} uma expressão regular sobre o nome do teste, com a
	justificativa escrita ao lado.
	\item \textbf{Arquivo:} para os arquivos temáticos sem identificador, como
	\texttt{test\_custodia.py}, o requisito do arquivo.
\end{enumerate}
Os níveis~2 e~3 estão num mapa explícito
(\texttt{docs/sistema/rastreabilidade-mapa.json}), escrito pelo autor com
apoio de IA, a partir do nome de cada teste. O \textit{script} para com erro
se o mapa citar um arquivo que não existe e avisa quando um identificador não
está no \textit{backlog}. Títulos genéricos, como ``horizonte~2'', não viram
histórias. Os \rastExcluidos{} testes do cliente simulado do \textit{front-end} ficam fora,
porque testam o simulador, e não o sistema.

% Gerado por docs/sistema/gera_rastreabilidade.py. Não editar à mão.
% Critério da ligação em docs/sistema/rastreabilidade-mapa.json; lista completa em docs/sistema/rastreabilidade.md.
\begin{scriptsize}
\begin{longtable}{>{\raggedright\arraybackslash}p{2.5cm}>{\raggedright\arraybackslash}p{1.9cm}rr>{\raggedright\arraybackslash}p{3.3cm}>{\raggedright\arraybackslash}p{4.5cm}}
\caption{Rastreabilidade entre requisitos e testes automatizados}\label{tab:rastreabilidade-testes}\\
\hline
\textbf{Requisito} & \textbf{Backlog} & \textbf{BE} & \textbf{FE} & \textbf{Arquivos (testes)} & \textbf{Exemplos} \\ \hline\endfirsthead
\caption[]{Rastreabilidade entre requisitos e testes automatizados (continuação)}\\
\hline
\textbf{Requisito} & \textbf{Backlog} & \textbf{BE} & \textbf{FE} & \textbf{Arquivos (testes)} & \textbf{Exemplos} \\ \hline\endhead
RF01 Coleta multifonte & feito: H01, H19; aguarda orientador: H50 & 62 & 18 & \texttt{AbaColeta.\allowbreak{}test.\allowbreak{}tsx} (17), \texttt{test\_\allowbreak{}conectores.\allowbreak{}py} (14), \texttt{test\_\allowbreak{}horizonte2.\allowbreak{}py} (14), \texttt{test\_\allowbreak{}editais\_\allowbreak{}pdf.\allowbreak{}py} (13), mais 6 & \texttt{test\_\allowbreak{}h19\_\allowbreak{}fabrica\_\allowbreak{}conhece\_\allowbreak{}os\_\allowbreak{}novos\_\allowbreak{}conectores}; \texttt{test\_\allowbreak{}h19\_\allowbreak{}portal\_\allowbreak{}envia\_\allowbreak{}a\_\allowbreak{}chave\_\allowbreak{}so\_\allowbreak{}no\_\allowbreak{}cabecalho} \\
RF02 Resolução de entidades & feito: H02 & 33 & 7 & \texttt{test\_\allowbreak{}cnpj\_\allowbreak{}alfanumerico.\allowbreak{}py} (11), \texttt{test\_\allowbreak{}resolucao.\allowbreak{}py} (8), \texttt{test\_\allowbreak{}normalizacao.\allowbreak{}py} (6), \texttt{cnpj.\allowbreak{}test.\allowbreak{}ts} (5), mais 5 & \texttt{test\_\allowbreak{}rf02\_\allowbreak{}socio\_\allowbreak{}servidor\_\allowbreak{}e\_\allowbreak{}doador\_\allowbreak{}nao\_\allowbreak{}sao\_\allowbreak{}fundidos\_\allowbreak{}e\_\allowbreak{}geram\_\allowbreak{}sugestao}; \texttt{test\_\allowbreak{}dv\_\allowbreak{}cnpj\_\allowbreak{}e\_\allowbreak{}valido} \\
RF03 Análise de vínculos em grafo & feito: H03, H40, H41, H42; parcial: H44, H57; fora desta entrega: H56 & 18 & 6 & \texttt{test\_\allowbreak{}lavagem.\allowbreak{}py} (12), \texttt{mapear.\allowbreak{}test.\allowbreak{}ts} (6), \texttt{test\_\allowbreak{}grafo.\allowbreak{}py} (5), \texttt{test\_\allowbreak{}api.\allowbreak{}py} (1) & \texttt{test\_\allowbreak{}rf03\_\allowbreak{}l4\_\allowbreak{}informa\_\allowbreak{}termino\_\allowbreak{}opaco\_\allowbreak{}e\_\allowbreak{}corta\_\allowbreak{}ciclo}; \texttt{test\_\allowbreak{}rf03\_\allowbreak{}l4\_\allowbreak{}dispara\_\allowbreak{}com\_\allowbreak{}dois\_\allowbreak{}niveis\_\allowbreak{}de\_\allowbreak{}socio\_\allowbreak{}pessoa\_\allowbreak{}juridica} \\
RF04 Alertas de tipologias & feito: H04, H31, H32, H33, H34, H35, H39, H40, H43, H47; parcial: H45, H49, H57; aguarda orientador: H46, H52, H53 & 72 & 34 & \texttt{test\_\allowbreak{}lavagem.\allowbreak{}py} (20), \texttt{test\_\allowbreak{}lavagem\_\allowbreak{}coleta.\allowbreak{}py} (17), \texttt{test\_\allowbreak{}horizonte2.\allowbreak{}py} (12), \texttt{AbaTipologias.\allowbreak{}test.\allowbreak{}tsx} (11), mais 9 & \texttt{test\_\allowbreak{}h31\_\allowbreak{}t4\_\allowbreak{}pela\_\allowbreak{}api\_\allowbreak{}e\_\allowbreak{}no\_\allowbreak{}relatorio}; \texttt{test\_\allowbreak{}rf04\_\allowbreak{}l16\_\allowbreak{}sinal\_\allowbreak{}b\_\allowbreak{}registro\_\allowbreak{}no\_\allowbreak{}cepim} \\
RF05 Linha do tempo patrimonial & aguarda orientador: H50, H51 & 0 & 0 & --- & nenhum teste \\
RF06 Monitoramento contínuo & feito: H18; parcial: H44 & 7 & 11 & \texttt{AbaMonitoramento.\allowbreak{}test.\allowbreak{}tsx} (11), \texttt{test\_\allowbreak{}horizonte2.\allowbreak{}py} (7) & \texttt{test\_\allowbreak{}h18\_\allowbreak{}agendador\_\allowbreak{}inicia\_\allowbreak{}e\_\allowbreak{}para\_\allowbreak{}com\_\allowbreak{}o\_\allowbreak{}servico}; \texttt{test\_\allowbreak{}h18\_\allowbreak{}fonte\_\allowbreak{}sem\_\allowbreak{}chave\_\allowbreak{}deixa\_\allowbreak{}a\_\allowbreak{}execucao\_\allowbreak{}parcial} \\
RF07 Cadeia de custódia digital & feito: H05, H14 & 15 & 6 & \texttt{AbaColeta.\allowbreak{}test.\allowbreak{}tsx} (6), \texttt{test\_\allowbreak{}horizonte1.\allowbreak{}py} (5), \texttt{test\_\allowbreak{}seguranca.\allowbreak{}py} (4), \texttt{test\_\allowbreak{}custodia.\allowbreak{}py} (3), mais 2 & \texttt{test\_\allowbreak{}h14\_\allowbreak{}recoleta\_\allowbreak{}do\_\allowbreak{}pncp\_\allowbreak{}compara\_\allowbreak{}por\_\allowbreak{}contrato}; \texttt{test\_\allowbreak{}h14\_\allowbreak{}comparacao\_\allowbreak{}e\_\allowbreak{}recoleta\_\allowbreak{}nao\_\allowbreak{}cruzam\_\allowbreak{}casos} \\
RF08 Gestão de casos & feito: H07, H37 & 4 & 15 & \texttt{AbaEquipe.\allowbreak{}test.\allowbreak{}tsx} (6), \texttt{Anotacoes.\allowbreak{}test.\allowbreak{}tsx} (5), \texttt{clienteHttp.\allowbreak{}h3.\allowbreak{}test.\allowbreak{}ts} (4), \texttt{test\_\allowbreak{}horizonte3.\allowbreak{}py} (3), mais 1 & \texttt{test\_\allowbreak{}h37\_\allowbreak{}tarefas\_\allowbreak{}com\_\allowbreak{}responsavel\_\allowbreak{}situacao\_\allowbreak{}prazo\_\allowbreak{}e\_\allowbreak{}historico}; \texttt{test\_\allowbreak{}h37\_\allowbreak{}anotacoes\_\allowbreak{}em\_\allowbreak{}entidade\_\allowbreak{}e\_\allowbreak{}alerta\_\allowbreak{}com\_\allowbreak{}autor\_\allowbreak{}data\_\allowbreak{}e\_\allowbreak{}auditoria} \\
RF09 Relatórios & feito: H08 & 8 & 5 & \texttt{test\_\allowbreak{}seguranca.\allowbreak{}py} (7), \texttt{AbaRelatorio.\allowbreak{}test.\allowbreak{}tsx} (5), \texttt{test\_\allowbreak{}api.\allowbreak{}py} (1) & \texttt{test\_\allowbreak{}relatorio\_\allowbreak{}json\_\allowbreak{}e\_\allowbreak{}markdown}; \texttt{test\_\allowbreak{}markdown\_\allowbreak{}escapa\_\allowbreak{}valores\_\allowbreak{}ja\_\allowbreak{}gravados} \\
RF10 Integrações sob ordem judicial & fora desta entrega: H25; parcial: H54; feito: H55 & 9 & 35 & \texttt{AbaOrdemJudicial.\allowbreak{}test.\allowbreak{}tsx} (16), \texttt{clienteMock.\allowbreak{}judicial.\allowbreak{}test.\allowbreak{}ts} (12), \texttt{test\_\allowbreak{}importacao\_\allowbreak{}judicial.\allowbreak{}py} (9), \texttt{clienteHttp.\allowbreak{}judicial.\allowbreak{}test.\allowbreak{}ts} (7) & \texttt{test\_\allowbreak{}rf10\_\allowbreak{}leiaute\_\allowbreak{}e\_\allowbreak{}o\_\allowbreak{}do\_\allowbreak{}anexo\_\allowbreak{}i\_\allowbreak{}da\_\allowbreak{}in\_\allowbreak{}bcb\_\allowbreak{}636}; \texttt{test\_\allowbreak{}rf10\_\allowbreak{}so\_\allowbreak{}dados\_\allowbreak{}sinteticos\_\allowbreak{}e\_\allowbreak{}pasta\_\allowbreak{}confinada\_\allowbreak{}a\_\allowbreak{}raiz} \\
RF11 Busca e filtros & feito: H09 & 4 & 0 & \texttt{test\_\allowbreak{}seguranca.\allowbreak{}py} (2), \texttt{test\_\allowbreak{}api.\allowbreak{}py} (1), \texttt{test\_\allowbreak{}cnpj\_\allowbreak{}alfanumerico.\allowbreak{}py} (1) & \texttt{test\_\allowbreak{}entradas\_\allowbreak{}de\_\allowbreak{}consulta\_\allowbreak{}tem\_\allowbreak{}limites}; \texttt{test\_\allowbreak{}busca\_\allowbreak{}trata\_\allowbreak{}curingas\_\allowbreak{}do\_\allowbreak{}like\_\allowbreak{}como\_\allowbreak{}texto} \\
RF12 Alerta como pista & feito: H12, H48 & 7 & 9 & \texttt{CartaoAlerta.\allowbreak{}test.\allowbreak{}tsx} (6), \texttt{test\_\allowbreak{}horizonte1.\allowbreak{}py} (6), \texttt{AbaAlertas.\allowbreak{}test.\allowbreak{}tsx} (3), \texttt{test\_\allowbreak{}combinacao.\allowbreak{}py} (1) & \texttt{test\_\allowbreak{}h12\_\allowbreak{}alerta\_\allowbreak{}de\_\allowbreak{}outro\_\allowbreak{}caso\_\allowbreak{}nao\_\allowbreak{}e\_\allowbreak{}decidido}; \texttt{test\_\allowbreak{}h12\_\allowbreak{}relatorio\_\allowbreak{}mostra\_\allowbreak{}a\_\allowbreak{}situacao\_\allowbreak{}de\_\allowbreak{}cada\_\allowbreak{}alerta} \\
RF13 Extração assistida por LLM & parcial: H22, H36 & 34 & 24 & \texttt{test\_\allowbreak{}horizonte3.\allowbreak{}py} (18), \texttt{AbaSugestoesIa.\allowbreak{}test.\allowbreak{}tsx} (15), \texttt{test\_\allowbreak{}editais\_\allowbreak{}pdf.\allowbreak{}py} (13), \texttt{localizarTrecho.\allowbreak{}test.\allowbreak{}ts} (4), mais 4 & \texttt{test\_\allowbreak{}h22\_\allowbreak{}blocos\_\allowbreak{}e\_\allowbreak{}provedor\_\allowbreak{}falso\_\allowbreak{}padrao}; \texttt{test\_\allowbreak{}h36\_\allowbreak{}t10\_\allowbreak{}cita\_\allowbreak{}o\_\allowbreak{}trecho\_\allowbreak{}literal\_\allowbreak{}e\_\allowbreak{}a\_\allowbreak{}posicao} \\
RF14 Priorização de alertas & feito: H13, H48 & 12 & 22 & \texttt{CartaoAlerta.\allowbreak{}test.\allowbreak{}tsx} (11), \texttt{AbaAlertas.\allowbreak{}test.\allowbreak{}tsx} (8), \texttt{test\_\allowbreak{}combinacao.\allowbreak{}py} (8), \texttt{AbaTipologias.\allowbreak{}test.\allowbreak{}tsx} (2), mais 3 & \texttt{test\_\allowbreak{}rf14\_\allowbreak{}pesos\_\allowbreak{}configurados\_\allowbreak{}sao\_\allowbreak{}validados}; \texttt{test\_\allowbreak{}h13\_\allowbreak{}tabela\_\allowbreak{}unica\_\allowbreak{}cobre\_\allowbreak{}todas\_\allowbreak{}as\_\allowbreak{}tipologias} \\
RF15 Análise temporal e geográfica & a fazer: H21 & 0 & 0 & --- & nenhum teste \\
RNF01 Segurança & feito: H10; parcial: H23 & 34 & 36 & \texttt{test\_\allowbreak{}horizonte3.\allowbreak{}py} (23), \texttt{Entrada.\allowbreak{}test.\allowbreak{}tsx} (15), \texttt{AbaEquipe.\allowbreak{}test.\allowbreak{}tsx} (10), \texttt{clienteHttp.\allowbreak{}h3.\allowbreak{}test.\allowbreak{}ts} (9), mais 4 & \texttt{test\_\allowbreak{}h23\_\allowbreak{}senha\_\allowbreak{}argon2id\_\allowbreak{}e\_\allowbreak{}politica}; \texttt{test\_\allowbreak{}h23\_\allowbreak{}totp\_\allowbreak{}janela\_\allowbreak{}uso\_\allowbreak{}unico\_\allowbreak{}e\_\allowbreak{}formato} \\
RNF02 Auditoria imutável & feito: H06, H55 & 35 & 36 & \texttt{test\_\allowbreak{}seguranca.\allowbreak{}py} (21), \texttt{clienteMock.\allowbreak{}judicial.\allowbreak{}test.\allowbreak{}ts} (10), \texttt{AbaOrdemJudicial.\allowbreak{}test.\allowbreak{}tsx} (7), \texttt{clienteHttp.\allowbreak{}test.\allowbreak{}ts} (7), mais 7 & \texttt{test\_\allowbreak{}rnf02\_\allowbreak{}cada\_\allowbreak{}acesso\_\allowbreak{}ao\_\allowbreak{}dado\_\allowbreak{}sigiloso\_\allowbreak{}entra\_\allowbreak{}na\_\allowbreak{}auditoria}; \texttt{test\_\allowbreak{}ancora\_\allowbreak{}so\_\allowbreak{}avanca} \\
RNF03 Anonimato operacional & fora desta entrega: H26 & 0 & 0 & --- & nenhum teste \\
RNF04 Conformidade legal & feito: H11, H17, H55 & 20 & 25 & \texttt{test\_\allowbreak{}privacidade\_\allowbreak{}cpf.\allowbreak{}py} (12), \texttt{privacidade.\allowbreak{}test.\allowbreak{}ts} (11), \texttt{test\_\allowbreak{}horizonte1.\allowbreak{}py} (4), \texttt{Casos.\allowbreak{}test.\allowbreak{}tsx} (3), mais 10 & \texttt{test\_\allowbreak{}h17\_\allowbreak{}o\_\allowbreak{}ultimo\_\allowbreak{}dia\_\allowbreak{}do\_\allowbreak{}prazo\_\allowbreak{}ainda\_\allowbreak{}vale}; \texttt{test\_\allowbreak{}h17\_\allowbreak{}caso\_\allowbreak{}exige\_\allowbreak{}finalidade\_\allowbreak{}e\_\allowbreak{}prazo\_\allowbreak{}de\_\allowbreak{}retencao} \\
RNF05 Explicabilidade & feito: H04; parcial: H22 & 27 & 18 & \texttt{test\_\allowbreak{}horizonte3.\allowbreak{}py} (18), \texttt{AbaSugestoesIa.\allowbreak{}test.\allowbreak{}tsx} (12), \texttt{test\_\allowbreak{}tipologias.\allowbreak{}py} (9), \texttt{CartaoAlerta.\allowbreak{}test.\allowbreak{}tsx} (3), mais 1 & \texttt{test\_\allowbreak{}h22\_\allowbreak{}blocos\_\allowbreak{}e\_\allowbreak{}provedor\_\allowbreak{}falso\_\allowbreak{}padrao}; \texttt{test\_\allowbreak{}h22\_\allowbreak{}desligada\_\allowbreak{}por\_\allowbreak{}padrao\_\allowbreak{}e\_\allowbreak{}nada\_\allowbreak{}e\_\allowbreak{}enviado} \\
RNF06 Desempenho e escala & feito: H29 & 3 & 0 & \texttt{test\_\allowbreak{}horizonte3.\allowbreak{}py} (3) & \texttt{test\_\allowbreak{}h29\_\allowbreak{}tetos\_\allowbreak{}e\_\allowbreak{}erros\_\allowbreak{}das\_\allowbreak{}consultas}; \texttt{test\_\allowbreak{}h29\_\allowbreak{}caminho\_\allowbreak{}em\_\allowbreak{}grafo\_\allowbreak{}sintetico\_\allowbreak{}maior} \\
RNF07 Disponibilidade e resiliência & sem história & 0 & 0 & --- & nenhum teste \\
RNF08 Usabilidade & sem história & 0 & 0 & --- & nenhum teste \\
RNF09 Interoperabilidade & a fazer: H24; parcial: H38 & 4 & 10 & \texttt{ExportarGrafo.\allowbreak{}test.\allowbreak{}tsx} (6), \texttt{clienteHttp.\allowbreak{}h3.\allowbreak{}test.\allowbreak{}ts} (4), \texttt{test\_\allowbreak{}horizonte3.\allowbreak{}py} (4) & \texttt{test\_\allowbreak{}h38\_\allowbreak{}zip\_\allowbreak{}com\_\allowbreak{}manifesto\_\allowbreak{}auditoria\_\allowbreak{}e\_\allowbreak{}permissoes}; \texttt{test\_\allowbreak{}h38\_\allowbreak{}graphml\_\allowbreak{}valido\_\allowbreak{}com\_\allowbreak{}custodia\_\allowbreak{}em\_\allowbreak{}nos\_\allowbreak{}e\_\allowbreak{}arestas} \\
RNF10 Manutenibilidade & feito: H01 & 14 & 0 & \texttt{test\_\allowbreak{}conectores.\allowbreak{}py} (14) & \texttt{test\_\allowbreak{}converter\_\allowbreak{}capital}; \texttt{test\_\allowbreak{}cnpj\_\allowbreak{}arquivo\_\allowbreak{}erros} \\
RNF11 DevSecOps & parcial: H16 & 0 & 0 & --- & nenhum teste \\
RNF12 Resiliência da coleta & parcial: H15 & 16 & 4 & \texttt{test\_\allowbreak{}seguranca.\allowbreak{}py} (7), \texttt{test\_\allowbreak{}horizonte1.\allowbreak{}py} (5), \texttt{AbaColeta.\allowbreak{}test.\allowbreak{}tsx} (4), \texttt{test\_\allowbreak{}conectores.\allowbreak{}py} (4) & \texttt{test\_\allowbreak{}h15\_\allowbreak{}retentativa\_\allowbreak{}comum\_\allowbreak{}conta\_\allowbreak{}tentativas\_\allowbreak{}e\_\allowbreak{}desiste}; \texttt{test\_\allowbreak{}h15\_\allowbreak{}cnpj\_\allowbreak{}arquivo\_\allowbreak{}tenta\_\allowbreak{}de\_\allowbreak{}novo\_\allowbreak{}em\_\allowbreak{}erro\_\allowbreak{}de\_\allowbreak{}leitura} \\
\hline
\multicolumn{6}{p{15.0cm}}{\scriptsize Fonte: gerada a partir dos repositórios do protótipo. BE: funções de teste do \textit{back-end} (um teste parametrizado conta uma vez); FE: blocos \texttt{it} do \textit{front-end}. Um teste pode ligar-se a mais de um requisito. Ligados: 360 de 390 no \textit{back-end} e 244 de 264 no \textit{front-end} (fora os 55 testes do cliente simulado). Critério em \texttt{docs/sistema/rastreabilidade-mapa.json}.}\\
\end{longtable}
\end{scriptsize}


A Tabela~\ref{tab:rastreabilidade-testes} mostra o resultado. Do
\textit{back-end}, \rastLigBack{} das \rastTestesBack{} funções de teste
ligam-se a algum requisito, \rastNomeBack{} delas pelo nome. Do
\textit{front-end}, fora os do cliente simulado, \rastLigFront{} dos
\rastElegiveisFront{} testes ligam-se, \rastNomeFront{} pelo título. No
\textit{back-end}, os testes sem ligação são sobretudo de migração do esquema,
de conferência das \textit{fixtures} e de limites defensivos dos conectores;
no \textit{front-end}, são sobretudo os do cliente HTTP, que conferem as rotas
e os cabeçalhos do contrato da API. A lista completa está em
\texttt{docs/sistema/rastreabilidade.md}. A contagem é de funções e de blocos
\texttt{it}: um teste parametrizado conta uma vez.
% PENDENTE: o autor revisa o mapa de rastreabilidade (níveis 2 e 3), feito com IA.

Seis requisitos não têm nenhum teste ligado. Três deles não foram
implementados: o RF05 aguarda decisão do orientador (H50 e H51), o RF15 está
por fazer (H21) e o RNF03 está fora desta entrega (H26). Os outros três estão
no sistema, mas sem teste automatizado:
\begin{itemize}
	\item \textbf{RNF07 (disponibilidade):} não tem história. Só há a política
	de reinício dos contêineres no \texttt{docker-compose.yml}, sem teste de
	\textit{backup} nem de recuperação;
	\item \textbf{RNF08 (usabilidade):} não tem teste automatizado; a
	avaliação com usuários prevista no Capítulo~\ref{cap:avaliacao} ainda não
	foi executada;
	\item \textbf{RNF11 (DevSecOps):} a parte feita é a configuração do
	\textit{pipeline}, com análise estática de segurança
	(Seção~\ref{sec:desenv-testes-ci}). Não é um teste, e não há DAST, SBOM
	nem \textit{pentest}.
\end{itemize}
Ter teste ligado também não prova que o requisito foi atendido por inteiro.
O RF01, por exemplo, lista fontes que o protótipo não coleta, como o DODF e
as juntas comerciais, e o RNF06 tem só testes do resultado e dos tetos das
consultas, sem medir o tempo.
A tabela mostra onde há teste, e não o quanto cada requisito foi verificado.

% ---------------------------------------------------------------------------
\subsection{Execução e cobertura}
\label{sec:desenv-testes-execucao}

Na última rodada registrada no diário, em 18/11/2026, depois das
tipologias de lavagem, do RF10 e da revisão da L6, o \textit{back-end} tinha 502 testes
coletados pelo \textit{pytest}, dos quais 8 são específicos de PostgreSQL.
Com SQLite, 494 passaram e os 8 de PostgreSQL foram pulados; com um
PostgreSQL descartável, os 502 passaram. A análise de estilo (ruff) e a
análise estática de segurança (Bandit) não apontaram problemas. O
\textit{front-end} tinha 323 testes, em 28 arquivos. A contagem do
\textit{back-end} na Tabela~\ref{tab:rastreabilidade-testes} difere desses
números porque conta funções, e não casos: um teste parametrizado é um
só.

A cobertura de código ainda não é medida. Sem ela, não se sabe que parte do
código a suíte executa \cite[cap.~5, seç.~4]{swebok4}.
% PENDENTE: preencher com os números de cobertura e dos testes E2E quando o dev-mvp concluir

% ---------------------------------------------------------------------------
\subsection{Portões da integração contínua}
\label{sec:desenv-testes-ci}

Os dois repositórios têm \textit{pipeline} no GitHub Actions, disparado a
cada \textit{push} e a cada \textit{pull request}. Cada etapa é um comando
que termina com erro quando falha, sem \texttt{continue-on-error}: qualquer
falha reprova a mudança. Os portões são estes:
\begin{itemize}
	\item \textbf{\textit{Back-end}:} instalação com as versões travadas
	(\texttt{uv sync -{}-locked}); análise de estilo (\texttt{ruff check}) e
	de formatação (\texttt{ruff format -{}-check}); análise estática de
	segurança (Bandit, sobre o código e os \textit{scripts}); testes com
	SQLite e com um PostgreSQL~16 descartável, criado pelo próprio
	\textit{job}, com o papel da aplicação; e uma etapa que roda de novo os
	testes com ``postgres'' no nome, em quatro arquivos, e reprova o
	\textit{pipeline} se algum deles for pulado.
	\item \textbf{\textit{Front-end}:} instalação com as versões travadas
	(\texttt{npm ci}); análise de estilo (ESLint); testes (Vitest);
	\textit{build} (\texttt{tsc} e Vite, sem o modo simulado); uma busca que
	reprova o \textit{build} se os tokens de demonstração tiverem entrado nele;
	e auditoria das dependências, que reprova vulnerabilidade alta ou crítica
	(\texttt{npm audit -{}-audit-level=high}).
\end{itemize}
Testes, análise de estilo e análise estática de segurança como portão são a
segunda prática candidata da RSL~2 para código gerado por IA (B2, TA2; B3,
TA2; Seção~\ref{sec:resultados-rsl2-praticas}). A auditoria de dependências
é decisão do autor. O \textit{pipeline} não roda
o teste contra as fontes reais nem a medição de desempenho, e ainda não mede
a cobertura.

Na \textit{branch} \texttt{develop}, em 17/11/2026, o \textit{pipeline} do
\textit{back-end} passou. No do \textit{front-end}, a análise de estilo, os
testes e o \textit{build} passaram, e a auditoria de dependências falhou: ela
aponta um aviso alto do pacote \texttt{picomatch} (GHSA-fc5q-39xc-vjpf), que
ainda não tem versão corrigida e só afeta dependências de desenvolvimento
(\texttt{npm audit -{}-omit=dev} não aponta nenhum).

% ---------------------------------------------------------------------------
\subsection{Teste com as fontes reais}
\label{sec:desenv-fontes-reais}

Em 17/11/2026, um \textit{script}
(\texttt{scripts/testar\_fontes\_reais.py}) rodou os conectores contra as
fontes reais, com as capturas só em memória, sem gravar no banco nem na
custódia. Houve só coleta passiva, com no máximo uma requisição por segundo
em cada fonte. Todas as fontes passaram: contratos e termos do PNCP, Portal
da Transparência (CEIS, CNEP e servidores, com a chave do operador), CNPJ
(amostra de 463~MB: três dos 30 ZIPs das tabelas lidas, de 2026-11, mais as
tabelas de domínio), TSE (candidatos e receitas do DF de
2026) e um edital real do PNCP em PDF, de 22 páginas, do qual foram
extraídos 63.398 caracteres em 0,43~s. Os dados reais mostraram quatro
divergências que os dublês não tinham revelado, corrigidas três no código e
uma na documentação:
\begin{itemize}
	\item o contrato do PNCP guardava o CPF completo do fornecedor pessoa
	física, e agora guarda só o descaracterizado;
	\item o sócio pessoa jurídica vem com o CNPJ básico, de 8 dígitos, e era
	tratado como pessoa; agora vira empresa;
	\item o arquivo de receitas do TSE de 2026 chama o ano de
	\texttt{AA\_ELEICAO}, e o conector passou a lê-lo;
	\item os limites de uso e o formato das datas do Portal da Transparência
	foram conferidos na fonte e registrados na documentação.
\end{itemize}
As três correções de código viraram testes de regressão, com dados
sintéticos que reproduzem o formato observado na fonte.
% Fontes: docs/mvp-diario.md (seção de 17/11/2026) e README do back-end
% (seção "Fontes reais e piloto com dados públicos"). O tempo do edital é o
% do diário (0,43 s); o README arredonda para 0,4 s.

% ---------------------------------------------------------------------------
\subsection{Desempenho do grafo}
\label{sec:desenv-desempenho}

A consulta de vínculos foi medida com grafos
sintéticos de até 1 milhão de arestas no PostgreSQL. Com 1 milhão de arestas
(500 mil nós), a vizinhança em até 3 níveis levou de 1,8 a 8,3~ms (mediana) e
o caminho entre duas entidades, 11,3~ms, contra 13 a 20 segundos da
implementação anterior, que montava o grafo inteiro em memória. O pico de
memória do processo caiu de 2.557~MiB para 83~MiB. Os números vêm de uma
execução por tamanho, numa máquina de desenvolvimento, e servem para comparar
ordens de grandeza (relatório completo em \texttt{docs/desempenho-grafo.md}
do \textit{back-end}).

% ---------------------------------------------------------------------------
\subsection{Lacunas e ameaças à validade dos testes}
\label{sec:desenv-testes-lacunas}

\begin{itemize}
	\item \textbf{Testes escritos pelo mesmo agente que escreveu o código.}
	A RSL~2 mostrou que testes gerados podem ter cobertura alta e, ainda
	assim, ser tautológicos ou ter asserções fracas (B2, TA2;
	Seção~\ref{sec:resultados-rsl2-qp2}). A prática candidata é o autor
	escrever ou validar os testes como critério de aceitação
	(Seção~\ref{sec:resultados-rsl2-praticas}). Aqui, o orquestrador rodou de
	novo cada suíte, o que confirma que os testes passam, mas não que eles
	são bons. A revisão humana dos testes está pendente.
	\item \textbf{Eficácia da suíte não medida.} Não há teste de mutação, que
	mede a parte dos defeitos injetados que a suíte detecta
	\cite[cap.~5, seç.~3.4 e~4.2.2]{swebok4}, nem cobertura de código. A
	RSL~2 aponta a mesma lacuna nos estudos: a eficácia dos testes gerados em
	achar defeitos raramente é medida (B2).
	\item \textbf{Dublês diferentes das fontes.} As \textit{fixtures} e os
	transportes simulados reproduzem o que se sabia do formato de cada fonte.
	O teste com as fontes reais achou quatro divergências que eles não
	cobriam, o que ilustra o risco descrito por Spadini et
	al.\ \cite{spadini2017mock}. Do mesmo modo, o cliente simulado do
	\textit{front-end} pode divergir da API real.
	\item \textbf{Sem navegador real.} Os testes do \textit{front-end} rodam
	no jsdom, que não é um navegador, e não há teste ponta a ponta. O grafo
	do Cytoscape.js, por exemplo, não é desenhado nos testes.
	\item \textbf{Dados aleatórios sem semente.} O sorteio dos CPFs e CNPJs a
	cada execução amplia as entradas, mas uma falha que dependa do número
	sorteado pode não se repetir.
	\item \textbf{Rastreabilidade por nome.} A ligação dos níveis~2 e~3 é
	decisão do autor e foi feita pelo nome dos testes. Um nome pode prometer
	mais do que o teste verifica.
\end{itemize}

% ===========================================================================
\section{Processo de desenvolvimento com IA}
\label{sec:desenv-processo}

O protótipo foi desenvolvido por agentes de codificação, sob orientação do
autor, e cada execução foi registrada num diário (\texttt{docs/mvp-diario.md}):
modelo e versão, tarefa, duração, tokens, resultado e correções. Foram nove
execuções. As de desenvolvimento usaram o modelo \texttt{claude-fable-5-1} e
depois o \texttt{claude-fable-5-2}; a orquestração, a revisão e a auditoria
de segurança usaram o \texttt{claude-opus-5-5}. Depois delas, as sessões de
16 e 17/11/2026 (teste da IA local, fontes reais e editais em PDF) foram
registradas em seções próprias do mesmo diário.

O processo aplicou as práticas candidatas da RSL~2
(Seção~\ref{sec:resultados-rsl2}):
\begin{itemize}
	\item \textbf{Verificação independente:} ao fim de cada execução, o
	orquestrador rodou de novo os testes, o lint, o \textit{build} e a análise
	de segurança, sem aceitar o relato do agente como prova;
	\item \textbf{Barreiras automáticas:} testes, análise estática e auditoria
	de dependências no \textit{pipeline};
	\item \textbf{Dados sintéticos:} só dados sintéticos nos testes
	automatizados e na demonstração. O teste de fumaça contra as fontes reais
	guarda as capturas só em memória, e os arquivos reais ficam fora do
	controle de versão;
	% PENDENTE: R1 registra se dados reais públicos chegaram ao contexto do
	% agente na sessão de 17/11/2026.
	\item \textbf{Registro de modelo, versão, tempo e tokens} de cada execução;
	\item \textbf{Fatos externos conferidos antes do código:} as regras das
	tipologias só usaram fatos conferidos em fonte oficial.
\end{itemize}
% PENDENTE: a RSL 2 aponta a revisão humana de cada mudança como prática
% central. Registrar como o autor revisou os diffs antes do commit.

% ===========================================================================
\section{Limitações e trabalho futuro}
\label{sec:desenv-limitacoes}

\begin{itemize}
	\item \textbf{Fontes reais:} as cinco fontes e um edital real em PDF
	passaram pelo teste de fumaça de 17/11/2026, e as quatro divergências
	encontradas foram corrigidas (Seção~\ref{sec:desenv-verificacao}). Foi um
	teste de conferência, e não uma investigação: não houve coleta para um
	caso, e a base completa do CNPJ não foi carregada. Essa base só entra por
	recorte, e o recorte não busca as outras empresas de um sócio do caso.
	Ficam pendentes: o CNPJ alfanumérico, que o código ainda trata como só
	dígitos; o CPF completo nas capturas brutas do PNCP, do CEIS e dos
	candidatos do TSE de 2026, e na comparação entre capturas; o arquivo de
	receitas do TSE de São Paulo (264~MB), acima do teto padrão de 64~MiB; e
	o nome do município do CNPJ, que não é resolvido.
	% PENDENTE: atualizar quando a correção do CPF/CNPJ alfanumérico terminar.
	\item \textbf{Integração:} a importação do GraphML e do CSV no Maltego e no
	i2 Analyst's Notebook não foi testada nos programas.
	\item \textbf{IA:} a extração foi executada uma vez com um servidor de
	modelo local real (Ollama 0.40.16, modelo \texttt{qwen2.5:7b}, temperatura
	0), sobre os três documentos sintéticos de editais da demonstração, em
	16/11/2026. O modelo devolveu 7 itens: 5 viraram sugestões e 2 foram
	recusados por trecho não literal. Um deles trazia o nome do órgão escrito
	errado pelo modelo, o que mostra a verificação do trecho barrando conteúdo
	inventado. O outro era um trecho verdadeiro que, no documento, estava
	quebrado em duas linhas, uma recusa indevida. Por isso a comparação passou
	a tratar qualquer sequência de espaços e quebras de linha como um espaço,
	mantendo exatas as letras, os acentos e a pontuação. Repetido o teste, o
	trecho foi aceito, e o nome escrito errado continuou recusado. Das 5
	sugestões, uma veio com
	o tipo errado (uma certidão classificada como objeto), e a decisão humana
	é que a corrige. Três documentos curtos não medem a qualidade da extração.
	Isso exige um conjunto anotado, que este trabalho não tem.
	\item \textbf{Chaves:} a troca da chave de criptografia não recifra as
	capturas antigas.
	\item \textbf{Editais em PDF:} só PDF com camada de texto é lido. Não há
	OCR, e o edital escaneado fica na custódia sem texto. O conector não
	baixa o edital do PNCP. O edital real foi só lido pelo extrator de texto,
	sem a extração por IA.
	\item \textbf{Avaliação com usuários:} o protótipo ainda não foi avaliado
	com policiais, como recomenda a síntese da QP5 da RSL~1.
	\item \textbf{Auditoria de dependências:} o passo de auditoria do
	\textit{pipeline} do \textit{front-end} falha por um aviso alto do
	\texttt{picomatch} (GHSA-fc5q-39xc-vjpf), sem versão corrigida, que só
	afeta dependências de desenvolvimento (Seção~\ref{sec:desenv-testes-ci}).
	\item \textbf{Tipologias de lavagem:} a leitura da tabela Simples do CNPJ,
	usada pela L6, não foi testada com o arquivo real, e a tela nova não foi
	conferida em navegador nem contra a API real
	(Seção~\ref{sec:desenv-lavagem}).
	\item \textbf{Dados abertos e lavagem de dinheiro:} com dados abertos, as
	tipologias de lavagem só apontam indícios cadastrais. Vários sinais de
	alerta de lavagem divulgados pelo Banco Central dependem de dado
	bancário, como depósitos em espécie relevantes em contas de servidores e
	de pessoas expostas politicamente e movimentações atípicas de recursos
	por agentes públicos \cite[art. 1º, I, n, e VII, a]{bcb2020cc4001}. Esse
	dado não é aberto. Por isso, o RF10 é uma importação separada, só para
	dado recebido por ordem judicial, no leiaute da Instrução Normativa BCB
	nº~636/2025 \cite{bcb2025in636}, e neste trabalho só aceita dados
	sintéticos (decisão do autor). O Relatório de Inteligência Financeira
	(RIF) do COAF chega por um sistema eletrônico de intercâmbio
	\cite[art. 1º]{coaf2026portaria60}, e não foi encontrado leiaute público
	dos dados dele.
	\item \textbf{O que o CNPJ aberto não traz:} os metadados dos dados
	abertos do CNPJ não têm o número de empregados, não dizem se o endereço é
	residencial ou comercial e trazem só o quadro atual de sócios, sem
	histórico \cite{rfbcnpjmetadados}. Assim, o sinal ``sem empregados
	registrados'', que o COAF associa a empresas recém-criadas com sócios de
	vínculo público \cite[tipologia 2.7]{coaf2021casos}, não é verificável.
	A troca de sócios só é vista entre duas coletas.
	\item \textbf{Parentesco:} nenhuma das fontes usadas traz vínculo
	familiar. O sobrenome igual não prova parentesco, e nenhuma regra o usa
	(decisão do autor).
	\item \textbf{Linha do tempo patrimonial (L11, RF05):} usa os bens que o
	candidato declara no registro de candidatura
	\cite[art. 11, § 1º, IV]{brasil1997lei9504}, publicados pelo TSE com o
	valor informado pelo próprio candidato \cite{tse2026dadosabertos}. Só
	cobre quem foi candidato, e a comparação exige ao menos duas eleições.
	Nos arquivos do Distrito Federal, o mesmo título eleitoral aparece em 299
	pessoas entre 2018 e 2022 e em 206 entre 2022 e 2026 (contagem do autor
	nos arquivos do TSE). A declaração de bens do agente público fica
	arquivada no serviço de pessoal e não é aberta
	\cite[art. 13]{brasil1992lei8429}. Renda privada, herança e renda do
	cônjuge não aparecem. E a desproporção, sozinha, é indício: a lei exige
	que os bens decorram de vantagem indevida e assegura ao agente demonstrar
	a licitude da origem \cite[art. 9º, VII]{brasil1992lei8429}.
	\item \textbf{Decisões pendentes do orientador:}
	\begin{itemize}
		\item \textbf{CPF dos candidatos.} Os arquivos de candidatos do TSE
		de 2018, 2022 e 2026 do Distrito Federal trazem o CPF completo, embora
		o leiame do TSE transcreva a regra de não divulgação do CPF e declare a
		variável não divulgável para as eleições de 2024
		\cite{tse2026dadosabertos}. Hoje o conector do TSE nunca extrai essa
		coluna para as entidades; o CPF completo que fica nas capturas brutas é
		a pendência do item ``Fontes reais'', acima. A
		proposta é lê-la só na ingestão e guardar apenas os 6 dígitos
		centrais, o mesmo recorte do CPF descaracterizado do CNPJ aberto
		\cite{rfbcnpjmetadados}, para comparar candidato, sócio e servidor
		como possível coincidência, sem fusão (decisão do autor, a validar
		pelo orientador). Isso muda a política atual, e dela dependem a L9 e o
		alerta da L11.
		\item \textbf{Remuneração de servidores distritais.} As rotas de
		servidores e de remuneração da API do Portal da Transparência são do
		Poder Executivo Federal \cite{cguapiportal}, e o download de servidores
		só traz servidores da União (Siape, Banco Central e militares)
		\cite{cguportalservidores}. A remuneração dos servidores do Governo do
		Distrito Federal (GDF) está no Portal da Transparência do Distrito
		Federal, em arquivos mensais com o CPF descaracterizado nos 6 dígitos
		centrais \cite{dfportalremuneracao}. Sem um conector para essa fonte, a
		T7, a L8 e a renda pública do RF05 não cobrem os servidores
		distritais.
	\end{itemize}
	% PENDENTE: decisão do orientador sobre o uso do CPF dos candidatos (6
	% dígitos centrais) e sobre o título eleitoral como chave entre eleições;
	% aceitar o Portal da Transparência do DF como fonte da remuneração distrital.
\end{itemize}

\endgroup
